\chapter{Formation Witnesses}
\label{ch:finite-depth-formation}

Exact parents are normalized pushforwards of the fine law
(\Cref{thm:rg-exact-coarse-vfe}) but in general not autonomous reduced
dynamics (\Cref{thm:rg-projection-memory}). This chapter asks whether, in one
normalized experiment at \(c_*\), coherent components can form in succession
while every parent is activated from the complete preceding recognition law,
retains every loop coordinate and source assignment, and learns its own copy
of a model parameter from its own law. Two witnesses answer yes under
explicit restrictions. \Cref{thm:finite-depth-adaptive-formation} treats a
triangle chain with a supplied hierarchy of interaction energies: at each
finite depth \(L\) the coherent components pass through
\(3\cdot2^{L-1}\to2^{L-1}\to\cdots\to1\) with probability tending to one,
while the discarded conditional losses, parameter covectors, assignment
histories, and continuous paths obey separate vanishing bounds.
\Cref{prop:rg-formation-endogenous-cycle} treats a homogeneous four-cycle
without supplied hierarchy, in which the source assignment decides which
pairs form, with probability bounded below independently of the small
parameter. \Cref{prop:rg-gated-depth-one} bounds what that mechanism can do
on any finite graph: in the gated family of
\Cref{def:rg-gated-translation-family}, the stationary coherent-block
partition is, up to the error of \cref{eq:rg-gated-block-law}, the
partition into the components of a minimum edge cover, and each merge
beyond such a cover costs a factor of the temperature. The proposition and
the conjecture of
the interpretation
chapter (\Cref{prop:phil-seeded-contrast-formation,conj:phil-contrast-formation})
are posed on the four-cycle family of the second witness. Both witnesses are
approximation statements under a stated preparation and regime, weaker than
the exact intertwining that an autonomous parent dynamics would require
(\Cref{open:rg-pointwise-parent-dynamics}). Both realize the general
construction for the translation group \((\R,+)\) on scalar sample fibers:
agents are local sections at \(c_*\) (\Cref{def:geo-agent}), one common frame
per agent lets one independent link translate both channels
(\Cref{def:geo-graph-links}), and the free energy is that of
\Cref{thm:obs-collective-vfe}.

Letters of the notation table that would collide with chapter-local objects
are replaced by the symbols of \cref{tab:rg-formation-notation}; upright
\(\mathrm J,\mathrm Q,\mathrm A,\mathrm T,\mathrm c_L,\mathrm C_L\) are local
scalars, the configuration is bold, and \(\ell'<\ell\) whenever two levels
occur.

\begin{sciencetable}{Chapter-local symbols of the formation families, shared with the proof appendix.}
\label{tab:rg-formation-notation}
\begin{tabularx}{\textwidth}{>{\raggedright\arraybackslash}p{0.28\textwidth}>{\raggedright\arraybackslash}X}
\toprule
\textbf{Symbol} & \textbf{Meaning} \\
\midrule
\(L\), \(N_\triangle=2^{L-1}\), \(N=3N_\triangle\), \(\ell,\ell'\) & depth, number of triangles, number of agents, level indices \\
\(\mathbf x=(\bar k,\bar m,\omega)\), \(\mathbf X_t\) & configuration of law presentations (belief means, model means, link translations) and its process \\
\(\mathfrak d_i=\bar k_i-\bar m_i\), \(\mathfrak D\), \(r_e\) & belief--model discrepancy, its row matrix, edge residual \\
\(\mathfrak h_q\), \(h_q\), \(O_\triangle\) & oriented triangle holonomy, normalized loop coordinate, contrast rows \\
\(\mathcal U,\mathcal H,\mathcal K\); \(W_\ell,\mathcal U_\ell,M_\ell,P_\ell\) & normal, loop, and gauge spaces; retained spaces and projections \\
\(H\), \(H_\ell\), \(\mathsf Q_\ell\) & stiffness, its Schur complement on \(M_\ell\), orthonormal flag columns \\
\(\mathrm J_\ell\), \(\mathrm Q_\ell\), \(\mathrm A\), \(\mathrm T\), \(t_*\) & energy strengths, their exponents, prepared amplitude, temperature, horizon \\
\(\sigma^2_{\mathrm{ev}}\), \(R_{\mathrm{ev}}\), \(\nu^*_i\) & evaluator variance, precision of the evaluator restriction cost, optimal evaluator coupling \\
\(\vartheta\), \(\Omega_0(\vartheta)\), \(\mathsf M\), \(\mathsf m=\mathrm T\mathsf M\) & learned scalar, prior precision, complete Fisher metric and its scaled form \\
\(\kappa\), \(z_\star\), \(o_{\mathcal K}\), \(o_{\mathrm{cal}}\), \(\mathsf o\) & root observation precision, root record, root and calibration observations, observed records \\
\(S\in\mathcal J\), \(\mathcal N_i^\triangle\), \(i_0\), \(\pi_{ij}(\vartheta)\), \(\pi_S(\vartheta)\) & complete source assignment, internal source palette, distinguished receiver, source priors \\
\(\chi\), \(\phi_S\), \(\delta\), \(B_0\), \(B_*\) & source function, source potential, record coefficient, and its two bounds \\
\(\beta_S\), \(\beta^\ell_{SS'}\), \(\bar\beta_{jk}\) & complete posterior source row, inherited proposal row, averaged fast row \\
\(q^\star_\vartheta\), \(q^{\star,\ell}_\vartheta\), \(\boldsymbol\Pi_{\vartheta,\mathsf o}\), \(q^\ell_t\) & restricted optimizer and its level-\(\ell\) push, physical posterior, recognition density \\
\(C_{\ell\leftarrow\ell'}\), \(\mathcal D^{\ell,\mu}_e\), \(\Delta_{\ell':\ell}\) & marked channel, inherited edge statistic, discarded conditional loss \\
\(\Fenergy_{\mathrm T}\), \(\mathrm C_{\mathrm T}\), \(\mathrm E_\vartheta\), \(U_{\vartheta,S}\), \(V_{\ell,S}\), \(\psi_{\ell,S}\), \(g_\ell\) & scaled restricted free energy, its offset, quadratic and sheet energies, parent potential and tilt, parameter covector \\
\(\mathscr L^\ell_\vartheta\) & inference generator at level \(\ell\) \\
\(t_\ell\), \(I_\ell\), \(t^{\mathrm{act}}_\ell\), \(k_L\), \(\eta_L\), \(\varsigma_\epsilon\), \(r_\epsilon\) & window start, window, activation time, entrance constant, tolerance, \(1+\log(1/\epsilon)\), parameter remainder \\
\(\mathrm c_L\), \(\mathrm C_L\), \(\mathsf j_m\) & positive constants depending on \(L\) and the fixed data; proposal marks \\
\(\mathsf B_{sr}\), \(\mathfrak j_\ell\), \(\mathfrak Y_\ell\), \(\mathfrak y_\ell\), \(\mathfrak x_\ell\), \(\varrho_s\), \(\mathrm Q_{\mathrm T}\), \(\mathrm Q_\Delta\) & flag blocks of the residual rows, least eigenvalue of a diagonal Gram block, forcing map and core coefficient of an interface, current-layer transient, strength ratio, exponents of the thermal ratio and of a loss target \\
\((V,E)\), \(\mathfrak n\), \(F\), \(\operatorname{cc}(F)\), \(\operatorname{att}(S)\) & graph of the gated family, its matching number, an edge set, the number of components of \((V,F)\), the edge set attended under \(S\) \\
\(\mathbf a_e\), \(\mathbf a_F\), \(W_F\), \(\mathrm s_F\) & residual pair map of an edge, its stack over \(F\), the kernel of the stack, the least positive singular value of the stack \\
\(H_S\), \(\mathrm k_S\), \(\mathrm h_S\), \(\Omega\), \(Z_S\), \(\mathrm z_S\), \(V_\pi\), \(V_{\min}\) & sheet stiffness, its rank and least positive eigenvalue, reference precision, sheet weight and its constant, source-marginal and gated energies \\
\(\operatorname{Coh}_{\mathrm a}(\mathbf x)\), \(\mathfrak P_{\mathrm a}(\mathbf x)\), \(\mathrm a\) & coherent edges and coherent blocks of a configuration; resolution \\
\bottomrule
\end{tabularx}
\end{sciencetable}

\paragraph{Scope.}
For the finite-depth family this chapter declines the following claims: that
attention or learning drives formation, since coherence is linear relaxation
of the prepared mean, the source record saturates at the coherence scale
(\Cref{rem:rg-formation-saturation}), and learning is calibration dominated up
to \(\epsilon^{2\mathrm Q_1+20}\); that per-agent models are slow, since means
and links have unit mobility and \(\mathfrak d_i\) is a level-one row; that
loop holonomy is of order one or nonabelian, since the group is \((\R,+)\) and
the retained loops are thermally negligible; that there is an inter-agent
model-channel record, since only the belief channel is observed; that parents
are coarse agents with their own belief and model fibers, since a parent
removes only normal coordinates; that source learning creates the supplied
contrasts; that the doubling recursion \(\mathrm Q_\ell=2\mathrm Q_{\ell+1}+20\)
is required, as distinct from the doubly exponential separation of inner base
\(4/3\) that \Cref{prop:rg-formation-gaussian-core} shows the Gaussian core to
require; that the continuous path laws converge in total variation or are exactly lumpable; and
that the construction saves computation or uses depth-uniform resources. For
the four-cycle it declines generic formation, formation of a matching from
every initial assignment, a unique block chosen by the invariant full law,
all-slow parents, exact closure under arbitrary source switching, permanent
membership, repeated coalescence, and a physical posterior conditioned on the
episode. The corresponding constructions are
\Cref{open:rg-formation-regime-sharpness,open:rg-formation-adaptive-family,open:rg-formation-general}.
\status{NOT-CLAIMED}

\section{The pointwise carrier and its restoring flag}
\label{sec:rg-formation-carrier}

\definitionheading{Triangle-chain carrier}{def:rg-formation-carrier}
Fix an integer \(L\ge2\), set \(N_\triangle=2^{L-1}\), and arrange
\(N_\triangle\) disjoint triangles in a chain, with one bridge between each
consecutive pair. Each triangle is cyclically oriented and every bridge has
one fixed orientation, so the graph has \(N=3N_\triangle\) vertices (the
agents) and \(4N_\triangle-1\) oriented edge copies. Bridge number \(j\),
between triangles \(j\) and \(j+1\), has level \(\ell(j)=2+v_2(j)\), where
\(v_2(j)\) is the exponent of two dividing \(j\); triangle edges have level
one. Energy strengths \(\mathrm J_1\ge\cdots\ge\mathrm J_L>0\) are attached to
the levels. The group \((\R,+)\) acts by translation on two scalar sample
fibers at \(c_*\); fix \(0<\sigma^2_{\mathrm{ev}}<1\). A configuration and its
law presentations are
\begin{equation}
\begin{aligned}
\mathbf x&=(\bar k,\bar m,\omega)\in\R^{2N+4N_\triangle-1}=\R^{10N_\triangle-1},&
\Theta_e(k)&=k+\omega_e,\\
q_i^b(\cdot\mid\mathbf x)&=\mathcal N(\bar k_i,1),&
q_i^m(\cdot\mid\mathbf x)&=\mathcal N(\bar m_i,\sigma^2_{\mathrm{ev}}),\\
\mathfrak d_i(\mathbf x)&=\bar k_i-\bar m_i,&
r_e(\mathbf x)&=\bar k_i-\bar k_j-\omega_e\quad(e:j\to i),
\end{aligned}
\label{eq:rg-formation-configuration}
\end{equation}
with \(\omega_{\bar e}=-\omega_e\). A passive frame change \(g_i\) sends
\(\bar k_i,\bar m_i,\omega_e\) to \(\bar k_i+g_i\), \(\bar m_i+g_i\),
\(\omega_e+g_i-g_j\), and re-expresses reference centers and observation
records in the same way. For triangle \(q\) put
\(\mathfrak h_q=\sum_{e\in q}\omega_e\) and \(h_q=\mathfrak h_q/\sqrt3\), and
let
\(O_\triangle=\begin{pmatrix}1/\sqrt2&-1/\sqrt2&0\\1/\sqrt6&1/\sqrt6&-2/\sqrt6\end{pmatrix}\).
Let \(B_1\) stack the \(N\) rows \(\mathfrak d_i\) and the \(2N_\triangle\)
internal contrast rows \(O_\triangle r_q\), where \(r_q\) is the vector of the
three residuals of triangle \(q\); for \(\ell\ge2\) let \(B_\ell\) stack the
bridge residuals of level \(\ell\), and let \(B\) stack all of them. With the
oriented incidence operator \((\partial t)_e=t_i-t_j\),
\begin{equation}
\begin{aligned}
\mathcal K&=\{(t,t,\partial t):t\in\R^N\},&
\mathcal H&=\operatorname{span}\{\text{normalized loop covectors}\},\\
\mathcal U&=\operatorname{range}B^\top,&
H&=B^\top\operatorname{diag}(\mathrm J_\ell I)B+\mathrm J_1P_{\mathcal H}.
\end{aligned}
\label{eq:rg-formation-orthogonal-split}
\end{equation}
The retained spaces and their Euclidean orthogonal projections are
\begin{equation}
W_\ell=\ker(B_1,\ldots,B_\ell),\qquad
\mathcal U_\ell=W_\ell\cap\mathcal U,\qquad
M_\ell=\mathcal U_\ell\oplus\mathcal H,\qquad
P_\ell:\R^{10N_\triangle-1}\to W_\ell,\qquad P_0=I.
\label{eq:rg-formation-projections}
\end{equation}

\lemmaheading{Carrier algebra}{lem:rg-formation-carrier}
In \Cref{def:rg-formation-carrier}, \(\mathfrak d_i\) and \(r_e\) are
invariant under passive frame changes, and every triangle satisfies
\begin{equation}
|r_q|^2=|O_\triangle r_q|^2+h_q^2,
\label{eq:rg-formation-loop-decomposition}
\end{equation}
so the loop energy belongs to the original sum of edge mismatches. The rows of
\(B\) are independent, \(\operatorname{rank}B=6N_\triangle-1\), the normalized
loop covectors are orthonormal and orthogonal to every row of \(B\), and
\(\R^{10N_\triangle-1}=\mathcal U\oplus\mathcal H\oplus\mathcal K\) is an
orthogonal direct sum. The stiffness has kernel \(\mathcal K\) and
\(\mathbf x^\top H\mathbf x=\mathrm J_1\sum_i\mathfrak d_i(\mathbf x)^2
+\sum_e\mathrm J_{\ell(e)}r_e(\mathbf x)^2\). Finally
\(P_\ell P_{\ell'}=P_\ell\) for \(\ell'<\ell\), every \(W_\ell\) contains
\(\mathcal H\oplus\mathcal K\), and
\begin{equation}
\dim W_\ell=4N_\triangle-1+\frac{N_\triangle}{2^{\ell-1}}\qquad(1\le\ell\le L).
\label{eq:rg-formation-dimensions}
\end{equation}
\status{ESTABLISHED}

\noindent\emph{Proof.}
A passive change adds \(g_i-g_i\) to \(\mathfrak d_i\) and
\(g_i-g_j-(g_i-g_j)\) to \(r_e\). Around an oriented triangle the vertex means
telescope, so \(\sum_{e\in q}r_e=-\mathfrak h_q\); the rows of \(O_\triangle\)
are orthonormal and orthogonal to \((1,1,1)/\sqrt3\), which gives
\cref{eq:rg-formation-loop-decomposition}. In a linear relation among the rows
of \(B\), the link columns occur only in contrast and bridge rows, with the
independent patterns \(-O_\triangle\) on the triangle edges and \(-1\) on each
bridge, so every contrast and bridge coefficient vanishes; the model-mean
columns then force the discrepancy coefficients to vanish. A normalized loop
covector has link entries \(1/\sqrt3\) on one triangle and nothing else, so it
is orthogonal to the discrepancy rows, to the contrast rows because
\(O_\triangle(1,1,1)^\top=0\), and to the bridges. Every row of \(B\) and every
loop covector annihilates \((t,t,\partial t)\), and
\((6N_\triangle-1)+N_\triangle+N=10N_\triangle-1\). The quadratic identity is
the sum of \(\mathrm J_\ell|B_\ell\mathbf x|^2\) and
\(\mathrm J_1|P_{\mathcal H}\mathbf x|^2\) combined with
\cref{eq:rg-formation-loop-decomposition}; it vanishes exactly when
\(B\mathbf x=0\) and \(P_{\mathcal H}\mathbf x=0\), that is on \(\mathcal K\).
The spaces \(W_\ell\) decrease, which gives \(P_\ell P_{\ell'}=P_\ell\), and
\(\mathcal H\oplus\mathcal K\) lies in each of them. There are
\(N_\triangle/2^{\ell''-1}\) bridges of level \(\ell''\ge2\), so
\(\operatorname{rank}(B_1,\ldots,B_\ell)=3N_\triangle+2N_\triangle
+N_\triangle(1-2^{1-\ell})\), which gives
\cref{eq:rg-formation-dimensions}. \(\square\)

The restoring estimates of the proof appendix rest on the bounds
\cref{eq:rg-formation-weighted-flag} of \(H\) by weighted energies
\(\sum_\ell\mathrm J_\ell|w_\ell|^2\) along orthonormal flag columns
\(\mathsf Q_\ell\) spanning \(\mathcal U_{\ell-1}\cap\mathcal U_\ell^\perp\).

\section{One normalized experiment with a nonlinear source record}
\label{sec:rg-formation-normalized-model}

\definitionheading{Normalized formation experiment}{def:rg-formation-experiment}
At each vertex the physical model sample is \(m_i\sim\mathcal N(\bar m_i,
\sigma^2_{\mathrm{ev}})\) and the evaluator sends \(m_i\) to
\(\mathcal N(m_i,1-\sigma^2_{\mathrm{ev}})\) on the belief sample \(k_i\).
Among couplings with the prescribed marginals
\(\mathcal N(\bar m_i,\sigma^2_{\mathrm{ev}})\) and \(\mathcal N(\bar k_i,1)\),
the recognition coupling at a fixed configuration is
\begin{equation}
m_i=\bar m_i+\xi_i,\qquad
k_i=\bar k_i+\xi_i+\xi'_i,\qquad
\xi_i\sim\mathcal N(0,\sigma^2_{\mathrm{ev}}),\quad
\xi'_i\sim\mathcal N(0,1-\sigma^2_{\mathrm{ev}}),
\label{eq:rg-formation-evaluator-lift}
\end{equation}
with independent noises; write \(\nu_i^*\) for its law and
\(R_{\mathrm{ev}}=\mathfrak D^\top\mathfrak D/(1-\sigma^2_{\mathrm{ev}})\),
where \(\mathfrak D\mathbf x=\bar k-\bar m\). Let \(\vartheta\in[1,21/10]\) be
the learned scalar and \(0<\mathrm T\le\mathrm J_1(1-\sigma^2_{\mathrm{ev}})/2\).
The configuration prior and its Lebesgue normalizer are
\begin{equation}
\begin{aligned}
\Omega_0(\vartheta)&=\vartheta H+P_{\mathcal K}-\mathrm TR_{\mathrm{ev}},\\
\mathbb P^0_\vartheta(d\mathbf x)
&=Z_0(\vartheta)^{-1}
  \exp\left[-\frac{\mathbf x^\top\Omega_0(\vartheta)\mathbf x}{2\mathrm T}\right]d\mathbf x,\\
Z_0(\vartheta)&=(2\pi\mathrm T)^{(10N_\triangle-1)/2}\det\Omega_0(\vartheta)^{-1/2}.
\end{aligned}
\label{eq:rg-formation-physical-prior}
\end{equation}
Fix \(0<\kappa\le1\) and a nonzero \(z_\star\in\mathcal K\). The observation
kernels are
\begin{equation}
o_{\mathcal K}\mid\mathbf x\sim
  \mathcal N_{\mathcal K}\bigl(P_{\mathcal K}\mathbf x,(\mathrm T/\kappa)I_{\mathcal K}\bigr),
\qquad
o_{\mathrm{cal}}\mid\vartheta\sim\mathcal N(\vartheta,\mathrm T),
\label{eq:rg-formation-observation-kernels}
\end{equation}
with the records \(o_{\mathcal K}=z_\star\) and \(o_{\mathrm{cal}}=2\); the
first is a density on the \(N\)-dimensional carrier \(\mathcal K\). Each
receiver \(i\) has the two candidate sources \(\mathcal N_i^\triangle\) inside
its own triangle, and the complete assignment is
\(S=(S_i)_i\in\mathcal J=\prod_i\mathcal N_i^\triangle\), with
\(|\mathcal J|=2^N\); every bridge is an added reversal pair of copies that keeps
its link and interaction energy but lies outside the source palettes. A distinguished receiver \(i_0\) has sources
labeled \(+\) and \(-\), and the independent source priors are
\begin{equation}
\begin{aligned}
\pi_{i_0,+}(\vartheta)
  &=\operatorname{sig}\bigl(4(\vartheta-3/2)\bigr),&
\pi_{i_0,-}(\vartheta)&=1-\pi_{i_0,+}(\vartheta),\\
\pi_{ij}(\vartheta)&=\tfrac12\quad(i\ne i_0),&
\pi_S(\vartheta)&=\prod_i\pi_{iS_i}(\vartheta),
\end{aligned}
\label{eq:rg-formation-source-priors}
\end{equation}
where \(\operatorname{sig}(u)=(1+e^{-u})^{-1}\). Set
\(B_0=10^{-4}\), \(\delta=B_0/N\), and \(\chi(r)=1-e^{-r^2/2}\). At each
receiver a binary record \(E_i\), conditionally independent given
\((\mathbf x,S)\), has
\begin{equation}
\begin{aligned}
\Pr(E_i=1\mid\mathbf x,S_i=j)
  &=q_0\exp\left[-\delta\chi\left(\sqrt{\mathrm J_1/\mathrm T} r_{ij}(\mathbf x)\right)\right],
       \qquad q_0=\tfrac12,\\
\phi_S(\mathbf x)&=\delta\sum_i
       \chi\left(\sqrt{\mathrm J_1/\mathrm T} r_{iS_i}(\mathbf x)\right),
       \qquad 0\le\phi_S\le B_0,
\end{aligned}
\label{eq:rg-formation-source-record}
\end{equation}
where \(r_{ij}\) is the residual of the edge from source \(j\) to receiver
\(i\), and the observed record is \(E_i=1\) for every \(i\). Since
\(\mathrm J_1r_{ij}^2/(2\mathrm T)=(\mathrm J_1/\mathrm T)
\KL(\mathcal N(\bar k_i,1)\Vert\mathcal N(\bar k_j+\omega_{ij},1))\), the
record is a bounded observation of the standardized edge mismatch. The
complete joint is
\begin{equation}
\begin{aligned}
\mathbb P_\vartheta
={}&\mathbb P^0_\vartheta(d\mathbf x)
 L_{\mathcal K}(do_{\mathcal K}\mid\mathbf x)
 L_{\mathrm{cal}}(do_{\mathrm{cal}}\mid\vartheta)\\
&{}\times\prod_i\left[
 \pi_i(dS_i\mid\vartheta)L_i(dE_i\mid\mathbf x,S_i)
 \mathcal N(dm_i;\bar m_i,\sigma^2_{\mathrm{ev}})
 \mathcal N(dk_i;m_i,1-\sigma^2_{\mathrm{ev}})\right].
\end{aligned}
\label{eq:rg-formation-full-joint}
\end{equation}

\propositionheading{Normalization and the evaluator optimum}{prop:rg-formation-normalized-experiment}
In \Cref{def:rg-formation-experiment}, for every admissible coupling
\(\nu_i\) of the two prescribed marginals,
\begin{equation}
\KL\bigl(\nu_i\Vert
  \mathcal N(\bar m_i,\sigma^2_{\mathrm{ev}})\operatorname{ev}_i\bigr)
=\KL(\nu_i\Vert\nu_i^*)+\frac{\mathfrak d_i^2}{2(1-\sigma^2_{\mathrm{ev}})},
\label{eq:rg-formation-evaluator-optimum}
\end{equation}
so \(\nu_i^*\) is the optimal coupling and the total restriction cost is
\(\mathbf x^\top R_{\mathrm{ev}}\mathbf x/2\). Moreover
\(\mathrm TR_{\mathrm{ev}}\preceq H/2\), hence
\(\Omega_0(\vartheta)\succeq(\vartheta/2)H+P_{\mathcal K}\succ0\) and the prior
is a normalized correlated Gaussian, and the joint
\cref{eq:rg-formation-full-joint} has total mass one, contains no kernel that
takes a recognition law as input, and gives the selected observation a
strictly positive density. \status{ESTABLISHED}

The term \(-\mathrm TR_{\mathrm{ev}}\) of the prior cancels the restriction
cost \(\tfrac{\mathrm T}2\mathbf x^\top R_{\mathrm{ev}}\mathbf x\) that
\cref{eq:rg-formation-evaluator-optimum} adds, so the restricted optimizer
below has a temperature-independent quadratic energy.

\section{Physical posterior, restricted optimizer, and recognition}
\label{sec:rg-formation-laws}

\definitionheading{Three laws of the formation experiment}{def:rg-formation-laws}
Write \(\mathsf o=(z_\star,2,\mathbf 1)\) for the observed records. The
recognition family consists of the laws with a smooth positive density
\(q(\mathbf x,S)\) relative to Lebesgue measure times counting measure and
conditional evaluator coupling \(\prod_i\nu_i^*\); assignments and
configurations may be correlated, and the conditional assignment law need not
be the source posterior. With
\begin{equation}
\begin{aligned}
\widehat K_\vartheta&=\vartheta H+(1+\kappa)P_{\mathcal K},&
\mathrm E_\vartheta(\mathbf x)&=\tfrac12\mathbf x^\top\widehat K_\vartheta\mathbf x
                  -\kappa z_\star^\top\mathbf x,\\
\mathrm E_\vartheta^{\mathrm{phys}}(\mathbf x)
   &=\mathrm E_\vartheta(\mathbf x)-\tfrac{\mathrm T}2\mathbf x^\top R_{\mathrm{ev}}\mathbf x,&
U_{\vartheta,S}(\mathbf x)&=\mathrm E_\vartheta(\mathbf x)+\mathrm T\phi_S(\mathbf x),\\
Z_{\mathrm{opt}}(\vartheta)
   &=\sum_S\pi_S(\vartheta)
      \int e^{-\mathrm E_\vartheta(\mathbf x)/\mathrm T-\phi_S(\mathbf x)}d\mathbf x,&
Z_{\mathrm{phys}}(\vartheta)
   &=\sum_S\pi_S(\vartheta)
      \int e^{-\mathrm E_\vartheta^{\mathrm{phys}}(\mathbf x)/\mathrm T-\phi_S(\mathbf x)}d\mathbf x,
\end{aligned}
\label{eq:rg-formation-partition-integrals}
\end{equation}
the restricted optimizer is
\(q^\star_\vartheta(\mathbf x,S)=\pi_S(\vartheta)e^{-U_{\vartheta,S}(\mathbf x)/\mathrm T}
/Z_{\mathrm{opt}}(\vartheta)\), and the physical posterior
\(\boldsymbol\Pi_{\vartheta,\mathsf o}\) on \((\mathbf x,S)\) is the same
expression with \(\mathrm E_\vartheta^{\mathrm{phys}}\) and
\(Z_{\mathrm{phys}}\). The scaled restricted free energy
\(\Fenergy_{\mathrm T}(q,\vartheta)\) is \(\mathrm T\) times the negative ELBO
of the evaluator restriction, evaluated at the records \(\mathsf o\).

\propositionheading{Laws, restricted free energy, and the physical gap}{prop:rg-formation-laws}
The integrals \cref{eq:rg-formation-partition-integrals} are finite and
positive, both laws have full support, and their common conditional
assignment row at fixed \(\mathbf x\) is the product of the receiver posterior
rows,
\begin{equation}
\beta_S(\mathbf x,\vartheta)=
\frac{\pi_S(\vartheta)e^{-\phi_S(\mathbf x)}}
     {\sum_{S'}\pi_{S'}(\vartheta)e^{-\phi_{S'}(\mathbf x)}}.
\label{eq:rg-formation-full-source-row}
\end{equation}
With the scalars \(\mathrm B_{\mathrm T}\), \(\mathrm C_{\mathrm T}\), and
\(p_\vartheta(\mathsf o)\) of \cref{eq:rg-formation-evidence-scalar},
\begin{equation}
\begin{aligned}
\Fenergy_{\mathrm T}(q,\vartheta)
 &=\mathrm T\sum_S\int q(\mathbf x,S)\log q(\mathbf x,S) d\mathbf x
  +\E_q\bigl[\mathrm E_\vartheta(\mathbf x)+\mathrm T\phi_S(\mathbf x)
       -\mathrm T\log\pi_S(\vartheta)\bigr]
       +\mathrm B_{\mathrm T}(\vartheta)\\
 &=\mathrm T\KL(q\Vert q^\star_\vartheta)+\mathrm C_{\mathrm T}(\vartheta),
\end{aligned}
\label{eq:rg-formation-restricted-vfe}
\end{equation}
and the excess over the unconstrained physical optimum is
\(\mathrm C_{\mathrm T}+\mathrm T\log p_\vartheta(\mathsf o)
=\mathrm T\log(Z_{\mathrm{phys}}/Z_{\mathrm{opt}})\ge0\). Whenever
\(N\mathrm T\le(1-\sigma^2_{\mathrm{ev}})\mathrm J_1\),
\begin{equation}
0\le\mathrm T\log\frac{Z_{\mathrm{phys}}(\vartheta)}{Z_{\mathrm{opt}}(\vartheta)}
\le\frac{2e^{B_0}N\mathrm T^2}{(1-\sigma^2_{\mathrm{ev}})\mathrm J_1},
\qquad
\KL(q^\star_\vartheta\Vert\boldsymbol\Pi_{\vartheta,\mathsf o})
\le\frac{2e^{B_0}N\mathrm T}{(1-\sigma^2_{\mathrm{ev}})\mathrm J_1}.
\label{eq:rg-formation-physical-optimizer-gap}
\end{equation}
\status{ESTABLISHED}

\section{Exact marked parents and their inherited inference}
\label{sec:rg-formation-parents}

\definitionheading{Marked parents and inherited edge statistics}{def:rg-formation-parents}
The parent carrier at level \(\ell\) is \(W_\ell\times\mathcal J\). For
\(\ell'<\ell\) the deterministic Markov kernel
\begin{equation}
C_{\ell\leftarrow\ell'}((y,S),\cdot)=\delta_{(P_\ell y,S)},
\qquad C_\ell=C_{\ell\leftarrow0},
\label{eq:rg-formation-marked-channel}
\end{equation}
extended by the identity on observations, pushes the generative joint, the
physical posterior, the restricted optimizer, and the recognition law, giving
\(q^{\star,\ell}_\vartheta=(C_\ell)_\#q^\star_\vartheta\),
\(\boldsymbol\Pi^\ell_{\vartheta,\mathsf o}\), and \((C_\ell)_\#q\), and
erases no assignment. The evaluator samples \((m_i,k_i)_i\), the records
\((E_i)_i\), and the discarded normal coordinates leave the parent carrier
through the disintegration of \cref{eq:rg-formation-full-joint} given
\((P_\ell\mathbf x,S)\), so their correlations are retained and no mark is
replaced by its value at \(P_\ell\mathbf x\). For an original edge
\(e\) and each \(\mu_\vartheta\in
\{\boldsymbol\Pi_{\vartheta,\mathsf o},q^\star_\vartheta\}\), the inherited
edge statistic is
\begin{equation}
\mathcal D_e^{\ell,\mu}(y,S)
=\E_{\mu_\vartheta}
 \left[\tfrac12r_e(\mathbf X)^2\mid P_\ell\mathbf X=y,S\right],
\label{eq:rg-formation-inherited-edge-statistic}
\end{equation}
the discrepancy observables for \(\mathfrak d_i\) are defined in the same way,
and at \(\ell=0\) both reduce to the configuration residuals. For a density
\(q^\ell\) on the marked parent carrier, the parent free energy is
\begin{equation}
\Fenergy^\ell_{\mathrm T}(q^\ell,\vartheta)
  =\mathrm T\KL(q^\ell\Vert q^{\star,\ell}_\vartheta)+\mathrm C_{\mathrm T}(\vartheta).
\label{eq:rg-formation-parent-vfe}
\end{equation}

\propositionheading{Exact parent potential and discarded conditional cost}{prop:rg-formation-parent-potential}
The kernels \cref{eq:rg-formation-marked-channel} compose,
\(C_{\ell''\leftarrow\ell}\circ C_{\ell\leftarrow\ell'}=C_{\ell''\leftarrow\ell'}\),
and positive smooth versions of every pushed law and conditional exist. The
root factor of \(q^{\star,\ell}_\vartheta\) is independent of the assignment
and of \(\vartheta\). On \(M_\ell\) let \(H_\ell\) be the Schur complement of
\(H|_{\mathcal K^\perp}\) and \(\gamma^{\mathrm{hid}}_{\ell,\vartheta}(dz\mid y)\)
the cumulative hidden Gaussian conditional of the quadratic optimizer, whose
mean is independent of \(\vartheta\) and whose covariance is
\(\mathrm T/\vartheta\) times a fixed inverse precision. Integration at fixed
\(S\) gives the sheet potential
\begin{equation}
\begin{aligned}
V_{\ell,S}(y,\vartheta)
 &=\tfrac12\vartheta y^\top H_\ell y+\mathrm C_\ell^{\mathrm G}(\vartheta)
       +\mathrm T\psi_{\ell,S}(y,\vartheta),\\
\psi_{\ell,S}(y,\vartheta)
 &=-\log\E_{\gamma^{\mathrm{hid}}_{\ell,\vartheta}(\cdot\mid y)}
                  e^{-\phi_S},&
0&\le\psi_{\ell,S}\le B_0,\\
\frac{d}{d\vartheta}\mathrm C_\ell^{\mathrm G}(\vartheta)
 &=\frac{\mathrm T \mathrm d_\ell}{2\vartheta},
\end{aligned}
\label{eq:rg-formation-parent-potential}
\end{equation}
where \(\mathrm d_\ell=\dim W_\ell^\perp\) is the cumulative eliminated
dimension, \(\mathrm C_0^{\mathrm G}=0\), \(\psi_{0,S}=\phi_S\), and the
unchanged root potential is appended on \(W_\ell\). At one common parameter
and one common time,
\begin{equation}
\begin{aligned}
\Delta_{\ell':\ell}(q^{\ell'},\vartheta)
 &=\E_{(C_{\ell\leftarrow\ell'})_\#q^{\ell'}}
   \KL\left(q^{\ell'}(dz\mid y,S)
        \middle\Vert q^{\star,\ell'}_\vartheta(dz\mid y,S)\right),\\
\Fenergy_{\mathrm T}^{\ell'}(q^{\ell'},\vartheta)
 &-\Fenergy_{\mathrm T}^\ell\bigl((C_{\ell\leftarrow\ell'})_\#q^{\ell'},\vartheta\bigr)
   =\mathrm T\Delta_{\ell':\ell}(q^{\ell'},\vartheta),
\end{aligned}
\label{eq:rg-formation-conditional-cost}
\end{equation}
where \(z\) consists only of the discarded normal coordinates, loops and
\(S\) staying in the retained state. \status{ESTABLISHED}

\noindent\emph{Proof.}
Composition follows from \(P_{\ell''}P_\ell=P_{\ell''}\), and finite Gaussian
domination supplies the positive smooth versions. The potential is the
Gaussian integral computed in \Cref{lem:formation-parent-score}. The cost
identity is \Cref{thm:rg-exact-coarse-vfe} for the channel
\(C_{\ell\leftarrow\ell'}\), since the scalar \(\mathrm C_{\mathrm T}\) is
common to both sides and the channel leaves the observations unchanged.
\(\square\)

\definitionheading{Declared inference kinetics}{def:rg-formation-kinetics}
Every microscopic coordinate has unit continuous mobility, and a rate-one
Poisson proposal clock redraws the whole assignment from
\(\beta(\mathbf x,\vartheta)\) without moving any continuous coordinate, a
proposal being allowed to return the current assignment. The microscopic
generator on global Euclidean space times the finite set \(\mathcal J\) is
\begin{equation}
\mathscr L_\vartheta^0\mathsf f(\mathbf x,S)
=\mathrm T\Delta_{\mathbf x}\mathsf f-\nabla_{\mathbf x}U_{\vartheta,S}\cdot\nabla_{\mathbf x}\mathsf f
 +\sum_{S'}\beta_{S'}(\mathbf x,\vartheta)[\mathsf f(\mathbf x,S')-\mathsf f(\mathbf x,S)],
\label{eq:rg-formation-microscopic-generator}
\end{equation}
with diffusion covariance \(2\mathrm TI dt\). The inherited proposal row and
parent generator are
\begin{equation}
\begin{aligned}
\beta^\ell_{SS'}(y,\vartheta)
 &=\E_{q^\star_\vartheta}
     [\beta_{S'}(\mathbf X,\vartheta)\mid P_\ell\mathbf X=y,S],\\
\mathscr L_\vartheta^\ell\mathsf f(y,S)
 &=\mathrm T\Delta_y\mathsf f-\nabla_yV_{\ell,S}\cdot\nabla_y\mathsf f
   +\sum_{S'}\beta^\ell_{SS'}(y,\vartheta)[\mathsf f(y,S')-\mathsf f(y,S)],
\end{aligned}
\label{eq:rg-formation-parent-generator}
\end{equation}
with the unchanged root drift included when \(y\) denotes all of \(W_\ell\).
Self-proposals are kept, so every row sums to one and every exit rate is at
most one; these rows declare the parent kinetics rather than project the fine
process in time.

\propositionheading{Reversible inherited rows}{prop:rg-formation-rate-composition}
The rows \(\beta^\ell\) satisfy detailed balance with \(q^{\star,\ell}_\vartheta\),
and for \(\ell'<\ell\)
\begin{equation}
\beta^{\ell'}_{SS'}(y,\vartheta)\ge e^{-B_0}\pi_{S'}(\vartheta)>0,
\qquad
\E_{q^{\star,\ell'}_\vartheta}
 [\beta^{\ell'}_{SS'}(Y,\vartheta)\mid P_\ell Y=z,S]
 =\beta^\ell_{SS'}(z,\vartheta).
\label{eq:rg-formation-rate-composition}
\end{equation}
\status{ESTABLISHED}

\noindent\emph{Proof.}
The product \(q^\star_\vartheta(\mathbf x,S)\beta_{S'}(\mathbf x,\vartheta)\)
equals the continuous marginal of \(q^\star_\vartheta\) times
\(\beta_S\beta_{S'}\) and is symmetric in \(S,S'\); conditional integration
over a retained fiber gives detailed balance. The lower bound follows from
\cref{eq:rg-formation-full-source-row} and \(0\le\phi_S\le B_0\), and the
second identity is the tower property of conditional expectation. \(\square\)

\section{Complete model scores and dissipation}
\label{sec:rg-formation-learning}

\definitionheading{Own-law parameter dynamics}{def:rg-formation-learning}
Let \(\mathsf M(\vartheta)\) be the Fisher information of the complete
normalized experiment \cref{eq:rg-formation-full-joint} in \(\vartheta\),
before conditioning on the records, and \(\mathsf m(\vartheta)=
\mathrm T\mathsf M(\vartheta)\). With the model metric
\(\mathsf m(\vartheta)/\epsilon\), \(\epsilon>0\), each level \(\ell\) evolves
its own recognition density and its own parameter by
\begin{equation}
\partial_tq_t^\ell=(\mathscr L_{\vartheta_\ell(t)}^\ell)^*q_t^\ell,
\qquad
\dot\vartheta_\ell(t)
  =-\epsilon \mathsf m(\vartheta_\ell(t))^{-1}
        g_\ell(q_t^\ell,\vartheta_\ell(t)),
\qquad
g_\ell=\left.\partial_\vartheta\Fenergy_{\mathrm T}^\ell\right|_{q^\ell}.
\label{eq:rg-formation-own-law-dynamics}
\end{equation}
The parameter path is deterministic, since the feedback is an expectation
under the law of its level. Parents inherit this metric, with no new rate
multiplier at activation.

\propositionheading{Complete metric, covector, and dissipation}{prop:rg-formation-score-dissipation}
The complete Fisher metric is
\begin{equation}
\mathsf M(\vartheta)
=\frac1{\mathrm T}+\frac12\Tr\bigl[(\Omega_0(\vartheta)^{-1}H)^2\bigr]
 +16\pi_{i_0,+}(\vartheta)\bigl[1-\pi_{i_0,+}(\vartheta)\bigr],
\label{eq:rg-formation-model-fisher}
\end{equation}
and the full parameter covector of level \(\ell\) is
\begin{equation}
\begin{aligned}
g_\ell(q^\ell,\vartheta)
={}&\vartheta-2+\tfrac12\E_{q^\ell}y^\top H_\ell y
       +\frac{\mathrm T \mathrm d_\ell}{2\vartheta}
       +\mathrm T\E_{q^\ell}\partial_\vartheta\psi_{\ell,S}(y,\vartheta)\\
 &-\frac{\mathrm T}2\Tr\bigl(\Omega_0(\vartheta)^{-1}H\bigr)
       -\mathrm T\E_{q^\ell}\partial_\vartheta\log\pi_S(\vartheta),
\end{aligned}
\label{eq:rg-formation-complete-score}
\end{equation}
a partial derivative at fixed recognition law, which at \(\ell=0\), where
\(\mathrm d_0=0\) and \(\partial_\vartheta\phi_S=0\), is the microscopic score.
Along \cref{eq:rg-formation-own-law-dynamics}, with
\(\mathfrak r_S=q^\ell_S/q^{\star,\ell}_{\vartheta_\ell}(\cdot,S)\),
\begin{equation}
\begin{aligned}
\frac{d}{dt}\Fenergy_{\mathrm T}^\ell
={}&-\mathrm T^2\sum_S\int q_S^\ell
       \left|\nabla\log\mathfrak r_S\right|^2dy
 -\frac{\mathrm T}2\int\sum_{S,S'}q^{\star,\ell}_{\vartheta_\ell}(y,S)\beta^\ell_{SS'}(y,\vartheta_\ell)
       (\mathfrak r_S-\mathfrak r_{S'})(\log\mathfrak r_S-\log\mathfrak r_{S'}) dy\\
 &-\epsilon \mathsf m(\vartheta_\ell)^{-1}g_\ell(q^\ell,\vartheta_\ell)^2
 \le0.
\end{aligned}
\label{eq:rg-formation-dissipation}
\end{equation}
At the microscopic level the jump term equals
\(-\mathrm T\int q_{\mathbf X}[\KL(q(\cdot\mid\mathbf x)\Vert\beta)
+\KL(\beta\Vert q(\cdot\mid\mathbf x))]d\mathbf x\). \status{ESTABLISHED}

The proofs are in \cref{sec:rg-formation-exact-proof}; the identities hold for
the solutions admitted by \Cref{lem:formation-prepared-continuation}, whose
cutoff passage is \Cref{lem:formation-activation-entropy}.

\section{Preparation and successive activation}
\label{sec:rg-formation-preparation}

\definitionheading{Regime, preparation, and activation schedule}{def:rg-formation-regime}
For \(0<\epsilon<1\) the regime is
\begin{equation}
\mathrm Q_L=0,\qquad \mathrm Q_\ell=2\mathrm Q_{\ell+1}+20,\qquad
\mathrm J_\ell=\epsilon^{-\mathrm Q_\ell},\qquad
\mathrm A^2=\epsilon^{3\mathrm Q_1+20},\qquad
\mathrm T=\mathrm A^2\epsilon^6,\qquad t_*=\frac2\epsilon,
\label{eq:rg-formation-weakest-one-regime}
\end{equation}
with unit mobility and proposal rate and the record coefficient \(\delta\) of
\cref{eq:rg-formation-source-record}; the weakest energy is exactly
\(\mathrm J_L=1\). Put \(\varsigma_\epsilon=1+\log(1/\epsilon)\) and
\begin{equation}
r_\epsilon=\mathrm A^2\mathrm J_1+\mathrm T\varsigma_\epsilon^2
=\epsilon^{2\mathrm Q_1+20}+\epsilon^{3\mathrm Q_1+26}\varsigma_\epsilon^2.
\label{eq:rg-formation-parameter-remainder}
\end{equation}
Let \(\mathbf v_\star\in\mathcal U\) be the unit vector of
\Cref{lem:formation-preparation-margin}, with margin \(\mathrm m_L>0\), and
\(\mathbf v_0\) any vector of the neighborhood of \(\mathbf v_\star\) on which
the margin is at least \(\mathrm m_L/2\). The initial parameter is
\(\vartheta(0)=1\). The root is independent of the non-root coordinates and
the assignments, with law \(\mathcal N_{\mathcal K}(\kappa z_\star/(1+\kappa),
\mathrm T(1+\kappa)^{-1}I_{\mathcal K})\). On every assignment sheet the
reference non-root law has common mean \(\mathrm A\mathbf v_0\), zero loop
mean, and covariance \(\mathrm T(H|_{\mathcal K^\perp})^{-1}\); the positive
assignment masses are near the product prior at \(\vartheta=1\), and the
distinguished source mass \(\lambda_{i_0,+}(0)=\Pr\{S_{i_0}(0)=+\}\) satisfies
\(|\lambda_{i_0,+}(0)-\operatorname{sig}(-2)|<10^{-2}\). The preparation
neighborhood consists of the smooth within-sheet transports of this reference
with bounded log-density ratio, retaining the root factor, whose coupling
satisfies
\begin{equation}
\bigl|\widehat K_1^{1/2}(\mathbf X_0-\mathbf X_0^{\mathrm G})\bigr|
 <\epsilon\sqrt{\mathrm T}
\quad\text{almost surely};
\label{eq:rg-formation-preparation-transport}
\end{equation}
small compactly supported diffeomorphisms in standardized non-root coordinates
supply such transports, the neighborhood may depend on \(\epsilon\) and on
\(L\), and later laws need not stay Gaussian or factor across sheets. Given an
entrance constant \(k_L\), the windows and activation times are
\begin{equation}
t_\ell=\frac{k_L\log(1/\epsilon)}{\mathrm J_\ell},
\qquad I_\ell=[t_\ell,2t_\ell],\qquad
t^{\mathrm{act}}_0=0,\quad t^{\mathrm{act}}_\ell=2t_\ell.
\label{eq:rg-formation-activation-times}
\end{equation}
At \(t^{\mathrm{act}}_\ell\) the complete level-\((\ell-1)\) law is pushed
through \(C_{\ell\leftarrow\ell-1}\), its parameter is copied, and the new law
then evolves by \cref{eq:rg-formation-own-law-dynamics}; an independently
continued source process keeps its own law and parameter.

\remarkheading{The regime in closed form}{rem:rg-formation-regime}
Solving the recursion gives \(\mathrm Q_\ell=20(2^{L-\ell}-1)\), so
\begin{equation}
\mathrm Q_1=20(2^{L-1}-1),\qquad
\frac{\mathrm J_1}{\mathrm J_L}=\epsilon^{-\mathrm Q_1},\qquad
\mathrm A^2=\epsilon^{3\mathrm Q_1+20},\qquad
\mathrm T=\epsilon^{3\mathrm Q_1+26}=\epsilon^{60\cdot2^{L-1}-34}.
\label{eq:rg-formation-regime-closed-form}
\end{equation}
The regime is a low-temperature limit with doubly exponential energy
separation in the depth: \(\mathrm T=\epsilon^{86},\epsilon^{206},
\epsilon^{446}\) for \(L=2,3,4\), every ratio \(\mathrm J_\ell/\mathrm T\)
diverges, \(\mathrm J_{\ell+1}/\mathrm J_\ell=\epsilon^{\mathrm Q_{\ell+1}+20}\)
and \(\mathrm J_{\ell+1}^2/\mathrm J_\ell=\epsilon^{20}\) for \(\ell<L\), and
the relative thermal scale is \(\sqrt{\mathrm T}/\mathrm A=\epsilon^3\). There
is no common acceleration of every energy, no vanishing source rate, and no
decrease of \(\delta\). \status{ESTABLISHED}

The recursion comes from the early estimate of
\Cref{lem:formation-conditional-entropy}, which bounds the retained Fisher
information by the full budget \cref{eq:rg-formation-full-fisher-budget} and
hence the discarded loss on \(I_\ell\) by the term
\((\mathrm A^2/\mathrm T)\mathrm J_{\ell+1}^2/\mathrm J_\ell\) of
\cref{eq:rg-formation-early-entropy-bound}. The linear-Gaussian core shows
that this term is not sharp.

\propositionheading{Gaussian core of the discarded conditional loss}{prop:rg-formation-gaussian-core}
Take \(\delta=0\), \(\vartheta\equiv1\), and an active non-root normal law
\(\mathcal N(\bar{\mathbf x}(t),\mathrm TH_{\ell'}^{-1})\) at level
\(\ell'<\ell\), whose covariance is stationary for the Ornstein--Uhlenbeck
flow, so \(\bar{\mathbf x}(t)=e^{-H_{\ell'}t}\bar{\mathbf x}(0)\). Split the
coordinates into the discarded \(z\) and retained \(y\), with blocks
\(H_{zz}\), \(H_{zy}\), \(H_{yy}\) of \(H_{\ell'}\) and Schur complement
\(H_\ell=H_{yy}-H_{zy}^\top H_{zz}^{-1}H_{zy}\).

(a) \emph{Loss formula and upper bound.} Then
\begin{equation}
\begin{aligned}
\Delta_{\ell':\ell}(t)&=\frac1{2\mathrm T}\mathfrak v^\top H_{zz}^{-1}\mathfrak v,
&
\mathfrak v&=(H_{\ell'}\bar{\mathbf x})_z,\\
\dot{\mathfrak v}&=-H_{zz}\mathfrak v-H_{zy}(H_{\ell'}\bar{\mathbf x})_y,
&
(H_{\ell'}\bar{\mathbf x})_y&=H_\ell\bar{\mathbf x}_y+H_{zy}^\top H_{zz}^{-1}\mathfrak v,
\end{aligned}
\label{eq:rg-formation-gaussian-core-loss}
\end{equation}
and, for \(|\bar{\mathbf x}|\le\mathrm C_L\mathrm A\),
\(\Delta_{\ell':\ell}\le\mathrm C_L(\mathrm A^2/\mathrm T)
[\mathrm J_{\ell'+1}e^{-\mathrm c_L\mathrm J_\ell(t-t^{\mathrm{act}}_{\ell'})}
+\mathrm J_{\ell+1}^4/\mathrm J_\ell^3]\).

(b) \emph{Two-sided asymptotic across an early interface.} Let the strengths
\(\mathrm J_1\ge\cdots\ge\mathrm J_L=1\) depend on \(\epsilon\) with ratios
\(\varrho_s=\mathrm J_{s+1}/\mathrm J_s\to0\) for \(s<L\), put \(\varrho_0=\varrho_L=0\), and
let \(1\le\ell\le L-1\) and \(\ell'=\ell-1\). Write
\(\mathsf B_{sr}=B_s\mathsf Q_r\) for the blocks of the matrix \(\mathsf B\)
of \cref{eq:rg-formation-weighted-flag}, which vanish for \(r>s\) and are
square and invertible for \(r=s\) (\Cref{lem:formation-parent-score}), let
\(\mathfrak j_\ell=\lambda_{\min}(\mathsf B_{\ell\ell}^\top\mathsf B_{\ell\ell})\),
let \(\mathfrak Y_\ell=\mathsf B_{\ell+1,\ell}^\top\mathsf B_{\ell+1,\ell+1}\mathsf B_{\ell+1,\ell+1}^\top\)
be the forcing map of the interface, and for
\(\mathbf v\in\R^{10N_\triangle-1}\) let
\begin{equation}
\mathfrak y_\ell(\mathbf v)=\tfrac12\bigl|(\mathsf B_{\ell\ell}^\top\mathsf B_{\ell\ell})^{-3/2}
\mathfrak Y_\ell B_{\ell+1}P_\ell\mathbf v\bigr|^2,
\qquad
\mathfrak x_\ell(t)=\varrho_\ell^{-2}\exp\bigl[-\tfrac12\mathfrak j_\ell\mathrm J_\ell(t-t^{\mathrm{act}}_{\ell-1})\bigr].
\label{eq:rg-formation-core-coefficient}
\end{equation}
Then \(\mathsf B_{\ell+1,\ell}\) has full row rank for every choice of the
bridge endpoints, so \(\mathfrak Y_\ell\) is injective. There are
\(\mathrm c_0,\mathrm C>0\) depending only on \(L\) and the carrier such
that, for the fine process started from the reference non-root law of
\Cref{def:rg-formation-regime} with mean \(\mathrm A\mathbf v_0\) and
activated at the times \cref{eq:rg-formation-activation-times}, whenever
\(\varrho_{\ell-1}+\varrho_\ell+\varrho_{\ell+1}\le\mathrm c_0\), the active level \(\ell-1\)
satisfies, for \(t^{\mathrm{act}}_{\ell-1}\le t\le t^{\mathrm{act}}_\ell\),
\begin{equation}
\begin{aligned}
\Bigl|\frac{\mathrm T\mathrm J_\ell^3}{\mathrm A^2\mathrm J_{\ell+1}^4}
\Delta_{\ell-1:\ell}(t)-\mathfrak y_\ell\bigl(\bar{\mathbf x}(t)/\mathrm A\bigr)\Bigr|
&\le\mathrm C\bigl[\varrho_{\ell-1}+\varrho_\ell+\varrho_{\ell+1}+\mathfrak x_\ell(t)+\mathfrak x_\ell(t)^2\bigr],\\
|B_{\ell+1}P_\ell\bar{\mathbf x}(t)-\mathrm AB_{\ell+1}P_\ell\mathbf v_0|
&\le\mathrm C\mathrm J_{\ell+1}\mathrm At.
\end{aligned}
\label{eq:rg-formation-core-two-sided}
\end{equation}
If moreover \(k_L\log(1/\epsilon)\varrho_\ell\to0\) and
\(\tfrac12\mathfrak j_\ell k_L\log(1/\epsilon)(1-2\varrho_{\ell-1})-2\log(1/\varrho_\ell)\to\infty\),
then, uniformly on \(I_\ell\),
\begin{equation}
\Delta_{\ell-1:\ell}(t)
=\frac{\mathrm A^2}{\mathrm T}\frac{\mathrm J_{\ell+1}^4}{\mathrm J_\ell^3}
\bigl[\mathfrak y_\ell(\mathbf v_0)+o(1)\bigr],
\qquad
\mathfrak y_\ell(\mathbf v_0)
\ge\frac{\sigma_{\min}(\mathfrak Y_\ell)^2N_\triangle\mathrm m_L^2}{2^{\ell+3}\|\mathsf B_{\ell\ell}\|^6}>0,
\label{eq:rg-formation-core-window}
\end{equation}
where \(\sigma_{\min}\) denotes the least singular value.

(c) \emph{Separation criterion.} With power laws
\(\mathrm J_s=\epsilon^{-\mathrm Q_s}\), \(\mathrm Q_1>\cdots>\mathrm Q_L=0\),
with \(\mathrm A^2/\mathrm T=\epsilon^{-\mathrm Q_{\mathrm T}}\) for some
\(\mathrm Q_{\mathrm T}>0\), which is \(\mathrm Q_{\mathrm T}=6\) in \cref{eq:rg-formation-weakest-one-regime}, and with an entrance
constant \(k_L>4(\mathrm Q_\ell-\mathrm Q_{\ell+1})/\mathfrak j_\ell\) for
every \(\ell<L\), both \(\inf_{I_\ell}\Delta_{\ell-1:\ell}\) and
\(\sup_{I_\ell}\Delta_{\ell-1:\ell}\) are of exact order
\(\epsilon^{3\mathrm Q_\ell-4\mathrm Q_{\ell+1}-\mathrm Q_{\mathrm T}}\). The core loss is
\(O(\epsilon^{\mathrm Q_\Delta})\) on every early window, \(\mathrm Q_\Delta\ge0\), exactly when
\(3\mathrm Q_\ell-4\mathrm Q_{\ell+1}\ge\mathrm Q_{\mathrm T}+\mathrm Q_\Delta\) for every \(\ell\le L-1\), which
forces \(\mathrm Q_1\ge(\mathrm Q_{\mathrm T}+\mathrm Q_\Delta)[(4/3)^{L-1}-1]\), with
equality along the recursion
\(3\mathrm Q_\ell=4\mathrm Q_{\ell+1}+\mathrm Q_{\mathrm T}+\mathrm Q_\Delta\); it vanishes on every
early window exactly when \(3\mathrm Q_\ell-4\mathrm Q_{\ell+1}>\mathrm Q_{\mathrm T}\) for every
\(\ell\le L-1\). \status{ESTABLISHED}

The proof of (b) and (c) is in \cref{sec:rg-formation-core-proof}.

\noindent\emph{Proof of (a).}
Both conditional laws of \(z\) given \(y\) have precision
\(H_{zz}/\mathrm T\), and their means differ by
\(H_{zz}^{-1}(H_{\ell'}\bar{\mathbf x})_z\) independently of \(y\), which gives
the first formula. Differentiating
\(\mathfrak v=(H_{\ell'}\bar{\mathbf x})_z\) along
\(\dot{\bar{\mathbf x}}=-H_{\ell'}\bar{\mathbf x}\) and eliminating
\(\bar{\mathbf x}_z\) give the other two. Put \(\mathfrak u=H_{zz}^{-1/2}\mathfrak v\). Its equation is
\[
\dot{\mathfrak u}
=-\bigl(H_{zz}+H_{zz}^{-1/2}H_{zy}H_{zy}^\top H_{zz}^{-1/2}\bigr)\mathfrak u
-H_{zz}^{-1/2}H_{zy}H_\ell\bar{\mathbf x}_y.
\]
The matrix multiplying \(\mathfrak u\) is symmetric and bounded below by
\(\mathrm c_L\mathrm J_\ell I\). The retained loop block is orthogonal to
\(H_{zy}\) and contributes no forcing; on the future normal coordinates,
\cref{eq:rg-formation-current-future-blocks} bounds the forcing by
\(\mathrm C_L\mathrm A\mathrm J_{\ell+1}^2/\sqrt{\mathrm J_\ell}\).
At \(t_0=t^{\mathrm{act}}_{\ell'}\), the weighted flag gives
\(|\mathfrak u(t_0)|\le\mathrm C_L\mathrm A\sqrt{\mathrm J_{\ell'+1}}\).
Variation of constants therefore gives
\[
|\mathfrak u(t)|\le\mathrm C_L\mathrm A\left[
\sqrt{\mathrm J_{\ell'+1}}e^{-\mathrm c_L\mathrm J_\ell(t-t_0)}
+\frac{\mathrm J_{\ell+1}^2}{\mathrm J_\ell^{3/2}}\right].
\]
Squaring and using \(\Delta_{\ell':\ell}=|\mathfrak u|^2/(2\mathrm T)\)
proves the bound, with the decay constant adjusted. The forcing vanishes for
\(\ell=L\). \(\square\)

\openproblemheading{Sharpness of the doubly exponential regime}{open:rg-formation-regime-sharpness}
Decide whether the recursion \(\mathrm Q_\ell=2\mathrm Q_{\ell+1}+20\) is
required for \Cref{thm:finite-depth-adaptive-formation}, and whether its
constants and preparation neighborhoods can be made uniform in the depth.
What holds is the theorem at every finite \(L\) with constants depending on
\(L\), and, by \Cref{prop:rg-formation-gaussian-core}, that in the Gaussian
core the discarded loss on every early window is of exact order
\((\mathrm A^2/\mathrm T)\mathrm J_{\ell+1}^4/\mathrm J_\ell^3\), smaller than
the term of the general estimate by \((\mathrm J_{\ell+1}/\mathrm J_\ell)^2\),
so that the early interfaces need exactly
\(3\mathrm Q_\ell\ge4\mathrm Q_{\ell+1}+20\) for the core loss at the scale
\(\epsilon^{14}\), a separation still doubly exponential in \(L\) but of inner
base \(4/3\), with \(\mathrm Q_\ell\ge14\) for the term
\(\mathrm C_{\mathcal J}/\mathrm J_\ell\) of the jump remainder
\cref{eq:rg-formation-conditional-jump-remainder}. The doubling enters the
present proof through the full Fisher budget
\cref{eq:rg-formation-full-fisher-budget} in the early estimate of
\Cref{lem:formation-conditional-entropy} and through the history estimate of
\Cref{lem:formation-history-coupling}, which uses it at the top interface in
\cref{eq:rg-formation-averaged-environment-error}. The depth enters through
\(\mathrm Q_1\) in the regime and every entrance condition, through
\(\log|\mathcal J|=N\log2\) in the jump remainder, and through the Frobenius
constants of \cref{eq:rg-formation-weighted-flag}, whose squares grow linearly
in \(N_\triangle\). The problem is posed in the family of the theorem under
(H1), (H3), and (H4), with power-law strengths
\(\mathrm J_\ell=\epsilon^{-\mathrm Q_\ell}\), \(\mathrm Q_L=0\), satisfying
\(3\mathrm Q_\ell\ge4\mathrm Q_{\ell+1}+20\) and \(\mathrm Q_\ell\ge14\) for
\(\ell<L\), with the amplitude, temperature, horizon, and record coefficient
of \cref{eq:rg-formation-weakest-one-regime} formed from \(\mathrm Q_1\), and
with an entrance constant above the thresholds of the theorem and of
\Cref{prop:rg-formation-gaussian-core}. It is settled by conclusions (a) to
(e) for every schedule of this class, with constants and preparation
neighborhoods independent of \(L\), through an early estimate of the retained
Fisher information by the future-layer energy, a history-environment bound
without the top-interface recursion, and depth-uniform flag and jump
constants; or by an admitted preparation and a schedule of the class under
which one of the conclusions fails, or a lower bound showing that one of the
conclusions cannot hold with constants independent of \(L\).

\section{The finite-depth formation theorem}
\label{sec:rg-formation-theorem}
\label{sec:rg-formation-common-regime-proof}

Formation in this family is hierarchical relaxation of the prepared mean. The
non-root configuration starts near \(\mathrm A\mathbf v_0\), where every
residual has order \(\mathrm A\) and every bridge residual
\(r_e(P_{\ell'}\mathrm A\mathbf v_0)\) with \(\ell'<\ell(e)\) stays at least
\(\mathrm A\mathrm m_L/2\) away from zero. The weighted flag makes residuals
of level at most \(\ell\) relax at rate of order \(\mathrm J_\ell\), while
bridges of higher level feel only couplings of order \(\mathrm J_{\ell+1}\).
On \(I_\ell\) the first have decayed by \(\epsilon^{\mathrm c_Lk_L}\) relative
to \(\mathrm A\), whereas the second have moved by at most
\(\mathrm C_L\mathrm A\mathrm J_{\ell+1}t_\ell\to0\) and keep their margin, so
each window joins pairs of coherent components along the bridges of level
\(\ell\) and halves their number. The source record reads residuals at the
thermal resolution \(\sqrt{\mathrm T/\mathrm J_1}\), far below
the coherence scale \(\mathrm A\), so it is saturated there. Source jumps, the
record, and learning are therefore perturbations, and the content of the
theorem is that exact marked parents, activation from the complete law, and
own-law learning can be carried through this relaxation with separately
vanishing errors.

\theoremheading{Finite-depth formation under a supplied energy hierarchy}{thm:finite-depth-adaptive-formation}
Assume
(H1) the experiment of \Cref{def:rg-formation-carrier,def:rg-formation-experiment,def:rg-formation-laws,def:rg-formation-parents}
with the kinetics of \Cref{def:rg-formation-kinetics,def:rg-formation-learning},
at a fixed depth \(L\ge2\);
(H2) the regime \cref{eq:rg-formation-weakest-one-regime};
(H3) an initial law in the preparation neighborhood
\cref{eq:rg-formation-preparation-transport};
(H4) parents activated at the times \cref{eq:rg-formation-activation-times}.
There exist constants \(k_{\mathrm{geom},L},\mathrm c_{\mathrm{early},L},
\mathrm c_{\mathrm{joint},L},\mathrm c_{\mathrm{fast},L},\eta_L^{\max}>0\),
depending only on \(L\) and the fixed data and constructed respectively in
\Cref{lem:formation-coherence} (together with \(\eta_L^{\max}\), a fixed
fraction of the margin \(\mathrm m_L\)),
\Cref{lem:formation-conditional-entropy},
\Cref{lem:formation-terminal-joint-entropy}, and
\Cref{lem:formation-fast-mixing}, with the following property. For every
entrance constant \(k_L\) with
\begin{equation}
\begin{aligned}
k_L&>k_{\mathrm{geom},L},&
\tfrac12\mathrm c_{\mathrm{early},L}k_L&>\mathrm Q_1+20,\\
\mathrm c_{\mathrm{joint},L}k_L&>\mathrm Q_1+8,&
\mathrm c_{\mathrm{fast},L}k_L&>\mathrm Q_1/2+12,
\end{aligned}
\label{eq:rg-formation-common-entrance}
\end{equation}
and every tolerance \(0<\eta_L\le\eta_L^{\max}\), there exist
\(\epsilon_L,\mathrm c_L,\mathrm C_L>0\), depending on \(L\), the fixed data,
\(k_L\), and \(\eta_L\), such that for every \(0<\epsilon<\epsilon_L\)
conclusions (a) to (e) hold for the fine process and for the lineage of
activated parents. The lineage is said to admit all \(L\) successive
activations when the fine process and every activated parent have a unique
nonexplosive solution of \cref{eq:rg-formation-own-law-dynamics} through
\(t_*\), with positive smooth densities and parameter in \([1,21/10]\).

(a) \emph{Well-posedness and parameter.} The lineage admits all \(L\)
successive activations, and every process satisfies, in absolute time,
\begin{equation}
1\le\vartheta_\ell(t)\le21/10,\qquad
|\dot\vartheta_\ell(t)|\le\mathrm C_L\epsilon,\qquad
|\vartheta_\ell(t)-(2-e^{-\epsilon t})|\le\mathrm C_Lr_\epsilon,
\label{eq:rg-formation-parameter-bound}
\end{equation}
with \(r_\epsilon=O(\epsilon^{2\mathrm Q_1+20})\) of
\cref{eq:rg-formation-parameter-remainder}; the comparison curve solves the
calibration part of \cref{eq:rg-formation-own-law-dynamics}.

(b) \emph{Coherence.} For a level-\(\ell\) description at time \(t\), an
original edge \(e\) is coherent when
\(\sqrt{2\mathcal D_e^{\ell,\mu}}\le\eta_L\mathrm A\), with
\cref{eq:rg-formation-inherited-edge-statistic} evaluated at the current
retained state and parameter of that description, and a vertex is internally
coherent when the corresponding observable for \(\mathfrak d_i\) is at most
\(\eta_L\mathrm A\); the coherent components are the connected components of
the graph on the \(N\) vertices whose edges are the coherent edges. For each
of \(\mu_\vartheta=\boldsymbol\Pi_{\vartheta,\mathsf o}\) and
\(\mu_\vartheta=q^\star_\vartheta\), except on an event of probability at most
\begin{equation}
\mathrm C_L(1+t_*)\exp(-\mathrm c_L/\epsilon^6),
\label{eq:rg-formation-coherence-probability}
\end{equation}
the microscopic process and its sequential parent descriptions have the
successive coherent components
\begin{equation}
3\cdot2^{L-1}\longrightarrow2^{L-1}
\longrightarrow2^{L-2}\longrightarrow\cdots\longrightarrow1.
\label{eq:rg-formation-agent-counts}
\end{equation}
On \(I_\ell\) every edge of level at most \(\ell\) is coherent, every vertex is
internally coherent, and every bridge of level greater than \(\ell\) has
residual observable at least a fixed positive multiple of \(\mathrm A\), so the
count on \(I_\ell\) is \(2^{L-\ell}\); the single final component persists
through \(t_*\). Loop coordinates remain random and retained, with no
identity-holonomy constraint.

(c) \emph{Conditional loss and score.} On each window \(I_\ell\), at the
current law and parameter,
\begin{equation}
\Delta_{\ell-1:\ell}
 \le
 \begin{cases}
\mathrm C_L\epsilon^{14},&\ell<L,\\
\mathrm C_L\epsilon^2,&\ell=L,
\end{cases}
\qquad
\bigl|g_{\ell-1}(q,\vartheta)
 -g_\ell\bigl((C_{\ell\leftarrow\ell-1})_\#q,\vartheta\bigr)\bigr|
 \le
 \begin{cases}
\mathrm C_L\mathrm T\epsilon^7,&\ell<L,\\
\mathrm C_L\mathrm T\epsilon,&\ell=L,
\end{cases},
\label{eq:rg-formation-conditional-score-bounds}
\end{equation}
and \(\mathrm T\Delta_{\ell-1:\ell}\) is the discarded restricted free energy.

(d) \emph{Source and kernel movement.} Let \(\lambda^\ell_{i_0,+}(t)\) be the
distinguished receiver's source mass in a fine, parent, or sequential process.
From each activation time at which that process is continued,
\begin{equation}
\lambda^\ell_{i_0,+}(t_*)-
  \lambda^\ell_{i_0,+}(t^{\mathrm{act}}_{\ell'})>\tfrac12,
\label{eq:rg-formation-source-movement}
\end{equation}
so the binary recognition source row changes by total variation greater than
one half. Between \(\vartheta=1\) and the final parameter, the first hidden
physical conditional and the first hidden optimizing conditional each change
by relative entropy greater than \(1/200\), at every fixed retained input and
assignment. Both changes are driven by the learned parameter through
\cref{eq:rg-formation-source-priors} and \(\vartheta H\), the record moving
the distinguished log odds by at most \(\delta\), and both take place after
formation, since \(\vartheta\le1+\mathrm C_L\epsilon\log(1/\epsilon)\) until
the last window closes.

(e) \emph{Histories and paths.} For any source process and its inherited
parent begun from a common marked activation sample at
\(t^{\mathrm{act}}_\ell\), there is a deterministic path
\(\bar{\mathbf x}_\ell:[t^{\mathrm{act}}_\ell,t_*]\to W_\ell\) with
\begin{equation}
\Pr\Bigl\{\sup_{t^{\mathrm{act}}_\ell\le t\le t_*}|P_\ell\mathbf X_t-\bar{\mathbf x}_\ell(t)|
  >\mathrm C_L\mathrm A\epsilon^2\Bigr\}
+\Pr\Bigl\{\sup_{t^{\mathrm{act}}_\ell\le t\le t_*}|\overline{\mathbf X}_t-\bar{\mathbf x}_\ell(t)|
  >\mathrm C_L\mathrm A\epsilon^2\Bigr\}
\le\mathrm C_L(1+t_*)e^{-\mathrm c_L/\epsilon^2}.
\label{eq:rg-formation-common-path-concentration}
\end{equation}
There is a coupling preserving the activation sample under which the first
bound below holds, and the second holds under every coupling of the two
processes:
\begin{equation}
\begin{aligned}
\Pr\bigl\{S_t\ne\overline S_t
       \text{ for some }t\in[t^{\mathrm{act}}_\ell,t_*]\bigr\}
 &\le\mathrm C_L\epsilon^6,\\
\Pr\Bigl\{\sup_{t^{\mathrm{act}}_\ell\le t\le t_*}
       |P_\ell\mathbf X_t-\overline{\mathbf X}_t|>\mathrm C_L\mathrm A\epsilon^2\Bigr\}
 &\le\mathrm C_L(1+t_*)e^{-\mathrm c_L/\epsilon^2}.
\end{aligned}
\label{eq:rg-formation-two-path-bounds}
\end{equation}
The resolution
\(\mathrm C_L\mathrm A\epsilon^2=\mathrm C_L\sqrt{\mathrm T}/\epsilon\) is
coarser than the thermal scale, and a finite number of successive
replacements changes only the constants. \status{ESTABLISHED}

\noindent\emph{Proof.}
The exact statements are
\Cref{prop:rg-formation-normalized-experiment,prop:rg-formation-laws,prop:rg-formation-parent-potential,prop:rg-formation-rate-composition},
and \Cref{tab:rg-formation-lemma-dependencies} lists the lemmas of
\cref{app:finite-depth-formation-proofs} with their inputs and exponents.
Conclusion (a) is
\Cref{lem:formation-prepared-continuation,lem:formation-parameter-comparison},
and (b) is \Cref{lem:formation-coherence}, which uses neither conditional
equilibration nor history averaging. For (c),
\Cref{lem:formation-conditional-entropy} gives the early bound from
\(\mathrm J_{\ell+1}^2/\mathrm J_\ell=\epsilon^{20}\) and the second entrance
condition, and \Cref{lem:formation-terminal-joint-entropy} the terminal bound
by joint diffusion-plus-jump relaxation; conditional relative entropy is
bounded by the full one through the chain rule, and the complete hidden score
bounds the covector difference by
\(\mathrm C_L\mathrm T(\sqrt\Delta+\Delta)\). Conclusion (d) is
\Cref{lem:formation-source-envelope}. For (e),
\Cref{lem:formation-history-coupling} couples the assignment histories and
\Cref{lem:formation-continuous-path} gives
\cref{eq:rg-formation-common-path-concentration}; each of its events concerns
the law of one process, so the union bound gives the second line of
\cref{eq:rg-formation-two-path-bounds} under every coupling. The four
entrance conditions are equivalent to the single threshold
\begin{equation}
k_L>\max\left\{k_{\mathrm{geom},L},\
\frac{2(\mathrm Q_1+20)}{\mathrm c_{\mathrm{early},L}},\
\frac{\mathrm Q_1+8}{\mathrm c_{\mathrm{joint},L}},\
\frac{\mathrm Q_1/2+12}{\mathrm c_{\mathrm{fast},L}}\right\},
\label{eq:rg-formation-joined-entrance-proof}
\end{equation}
whose terms supply, in order, the geometric windows, the early entrance,
terminal joint relaxation, and the source-history estimate; every lemma uses
the same preparation, regime, and entrance choice. The final window has order
\(k_L\log(1/\epsilon)\) and lies before \(t_*-10=2/\epsilon-10\) once
\(\epsilon\) is small; this, the ordering of the windows, and the smallness
conditions inside the lemmas fix \(\epsilon_L\) after \(k_L\) and \(\eta_L\).
\(\square\)

\remarkheading{Saturated record, calibrated learning, and thermally negligible loops}{rem:rg-formation-saturation}
Uniformly in the configuration, the assignment, and the parameter,
\begin{equation}
\begin{aligned}
|\log\beta_S(\mathbf x,\vartheta)-\log\pi_S(\vartheta)|&\le B_0=10^{-4},\\
|H^{-1/2}\mathrm T\nabla_{\mathbf x}\phi_S(\mathbf x)|&\le B_*\sqrt{\mathrm T},
\qquad
B_*=\delta\sqrt{2N}=10^{-4}\sqrt{2/N}\le5.8\times10^{-5},
\end{aligned}
\label{eq:rg-formation-source-saturation}
\end{equation}
with \(H^{-1/2}\) acting on \(\mathcal K^\perp\), by
\cref{eq:rg-formation-full-source-row} with \(0\le\phi_S\le B_0\) and by
\cref{eq:rg-formation-original-relative-bounds}. At a residual of order
\(\mathrm A\) the argument of \(\chi\) has order
\(\epsilon^{-(3+\mathrm Q_1/2)}\), so the record equals its saturated value up
to an exponentially small error. The proof of \Cref{lem:formation-coherence}
uses the record and the assignments only through \(0\le\phi_S\le B_0\) and the
pathwise force bound \cref{eq:rg-formation-path-force-bound}, valid for every
adapted assignment path, and the parameter only through
\(1\le\vartheta\le21/10\), so with \(\delta=0\) and \(\vartheta\equiv1\) it
gives the same components. By
\Cref{lem:formation-prepared-continuation,lem:formation-parameter-comparison}
every law-dependent term of the covector is at most \(\mathrm C_Lr_\epsilon\),
which gives the parameter bounds of \Cref{thm:finite-depth-adaptive-formation}.
The loop energy \(\mathrm J_1|h|^2\) belongs to the strongest term, so
\(\mathrm J_1\E|h|^2\le\mathrm C_L\mathrm T\) by
\cref{eq:rg-formation-actual-moments}, and the loops have root-mean-square size
at most \(\mathrm C_L\sqrt{\mathrm T/\mathrm J_1}
=\mathrm C_L\mathrm A\epsilon^{3+\mathrm Q_1/2}\). \status{ESTABLISHED}

\section{A two-triangle instance and the remaining scope}
\label{sec:rg-formation-example-scope}

\remarkheading{The two-triangle instance}{rem:rg-formation-two-triangle}
For \(L=2\), take triangles \((1,2,3)\) and \((4,5,6)\) with the single bridge
\((3,4)\). The retained carriers are
\begin{equation}
\R^{19}\times\mathcal J
\longrightarrow\R^9\times\mathcal J
\longrightarrow\R^8\times\mathcal J,
\qquad |\mathcal J|=64,
\label{eq:rg-formation-two-triangle-carrier}
\end{equation}
the coherent component counts are \(6\to2\to1\), both loops survive both
maps, and the regime becomes
\begin{equation}
\mathrm J_1=\epsilon^{-20},\qquad \mathrm J_2=1,\qquad
\mathrm A=\epsilon^{40},\qquad \mathrm T=\epsilon^{86}.
\label{eq:rg-formation-two-triangle-regime}
\end{equation}
The first window has order \(\epsilon^{20}\log(1/\epsilon)\), the second
\(\log(1/\epsilon)\), and the source and model changes occur over the longer
horizon \(2/\epsilon\). The source-dependent term is not quadratic: along a
direction with selected internal residual slope \(\lambda\ne0\),
\(\chi(\lambda t)=\lambda^2t^2/2-\lambda^4t^4/8+O(t^6)\), and every assignment
selects an internal edge with nonzero first-normal projection, so the hidden
log density has a nonzero fourth derivative in some first-normal direction;
exact marginalization keeps the full conditional integral and the negative
covariance term of its Hessian, and neither first hidden conditional is a
Gaussian. \status{ESTABLISHED}

\openproblemheading{A family in which learning or holonomy selects the components}{open:rg-formation-adaptive-family}
Construct a normalized family at \(c_*\) in which parameter learning or loop
holonomy changes, at the level of the law, which components form, with exact
marked parents and activation from the complete law. What holds is
\Cref{thm:finite-depth-adaptive-formation}, whose components are fixed by the
supplied contrasts and internal source support, and
\Cref{prop:rg-formation-endogenous-cycle}, whose pairs are seeded by the
random initial assignment without learning, so that in it the source record
selects only along sampled paths; \Cref{prop:phil-seeded-contrast-formation}
shows that its seeded bound is unchanged by learning and that from an initial
law invariant under the cycle automorphisms learning leaves the rows uniform,
so a law-level selection needs an initial law without that invariance. Such a
family needs at least one of a source record standardized at the coherence
scale, for example through \(\chi(r/(\eta_L\mathrm A))\), model means and
model links both slow relative to the belief mobility, since a fast link or
fast model means remove the model disagreement on the belief clock whatever
the other rate (\Cref{prop:obs-slow-model-link-dichotomy}), loop coordinates
of order \(\mathrm A\) or nonabelian links, or interacting model-channel
records, together with new fast-marginal and history estimates; sources on
different bridge levels and source learning that creates the contrasts belong
to the same problem. It is posed on the four-cycle of
\Cref{prop:phil-seeded-contrast-formation}, with its temperature, learning
mobility, and horizon, for Gaussian initial laws
\(\mathcal N(\mathsf m_0,I_{12})\) with \(\mathsf m_0\) outside a
Lebesgue-null set, and for parents of the type of part (b) of
\Cref{prop:rg-formation-endogenous-cycle} activated from the recognition law
of the episode prefix; the law of the formed block means the conditional law,
given contrast formation, of the first block for which it occurs, the
conditioning event having probability bounded below in every arm by part (b)
of \Cref{prop:phil-seeded-contrast-formation}, so that a difference in
whether formation occurs is not selection. It is settled by such parents with
separately vanishing conditional and dynamical errors together with a
total-variation difference, bounded below as \(\epsilon\to0\), between the
conditional laws of the formed block in the learned and in the frozen-row
arm, or between those of two initial means that differ only in their loop
coordinate \(\sum_i\omega_i\), the amplified datum being the initial law or
its loop coordinate and the other arm its null; or by a proof that, for
almost every \(\mathsf m_0\), each of these two pairs of arms gives
conditional laws of the formed block whose difference vanishes as
\(\epsilon\to0\).

A separate slower-mobility regime is resolved by
\cref{prop:phil-early-block-null}: keeping the same contrast event and
horizon but taking \(\epsilon_{\mathrm L}=o(\mathrm T^{1/4})\), the
learned and frozen conditional first-block laws have vanishing total
variation for every fixed Gaussian mean. The original coupled mobility
\(\epsilon_{\mathrm L}=\mathrm T^{1/8}\) and the loop-coordinate
comparison remain outside that estimate. The late coherent-partition
contrast in \cref{prop:phil-selected-pair-persistence} concerns restored
equilibrium coherence and does not close this first-block problem.
For the late coherence event on the joint slow schedule,
\cref{prop:phil-actual-conditioned-parent,prop:phil-actual-parent-dynamics}
establish actual conditioned activation of the existing parent with
separately vanishing conditional-information and dynamical errors.
The activation event is a completed coherence prefix after selection.
It is not the first-contrast event in the original regime posed above.
For the seeded sub-event of that original contrast, at every mobility
\(\epsilon_{\mathrm L}\in[0,1]\) including \(\mathrm T^{1/8}\),
\cref{prop:phil-early-seeded-parent,prop:phil-early-parent-dynamics} give
separately vanishing conditional and dynamical errors when the parent is
activated after hidden thermalization, at a deterministic activation time
\(\theta\) with \(\theta-\log(1/\epsilon)\to\infty\) and
\(\epsilon^4\theta\to0\), on following windows short compared with
\(\mathrm T^{-1}\), while
\cref{prop:phil-early-first-completion} shows that activation at the first
qualifying completion has diverging conditional error, also under
conditioning on the whole contrast event. The learned and frozen first-block
comparison at \(\epsilon_{\mathrm L}=\mathrm T^{1/8}\), the loop-coordinate
comparison, and conditioning on contrast histories outside the seeds remain
open.

\section{An endogenous four-cycle witness}
\label{sec:rg-formation-endogenous-cycle}

\definitionheading{Homogeneous four-cycle experiment}{def:rg-formation-cycle-experiment}
Index four agents by \(i\in\mathbb Z/4\) and orient edge \(e_i\) from \(i\) to
\(i+1\). A configuration is \(\mathbf x=(\bar k,\bar m,\omega)\in\R^{12}\), with
\(q_i^b=\mathcal N(\bar k_i,1)\), \(q_i^m=\mathcal N(\bar m_i,\sigma^2_{\mathrm{ev}})\),
\(\sigma^2_{\mathrm{ev}}=1/2\), four independent link translations
\(\omega_i\), residuals \(r^b_i=\bar k_{i+1}-\bar k_i-\omega_i\) and
\(r^m_i=\bar m_{i+1}-\bar m_i-\omega_i\), and the edge cost
\begin{equation}
\mathcal D_{e_i}(\mathbf x)=\tfrac12\bigl[(r^b_i)^2+(r^m_i)^2\bigr]
=\KL\bigl(q^b_{i+1}\Vert(\Theta_{e_i})_\#q^b_i\bigr)
+\sigma^2_{\mathrm{ev}}\KL\bigl(q^m_{i+1}\Vert(\Theta_{e_i})_\#q^m_i\bigr),
\label{eq:rg-formation-cycle-cost}
\end{equation}
which is the same for both orientations. Every receiver has both neighbors as
sources, \(S\in\mathcal J=\prod_i\{i-1,i+1\}\) with \(|\mathcal J|=16\) and
uniform prior \(\pi_S=1/16\), and \(V_S=\sum_i\mathcal D_{e(i,S_i)}\), where
\(e(i,j)\) joins \(i\) and \(j\). The configuration prior is
\(\mathcal N(0,I_{12})\); given \((\mathbf x,S)\), independent binary records
have \(\Pr(E_i=1\mid\mathbf x,S)=\tfrac12\exp[-\mathcal D_{e(i,S_i)}(\mathbf x)/\mathrm T]\),
and \(E=\mathbf 1\) is observed; the physical evaluator and recognition
coupling are those of \Cref{def:rg-formation-experiment}. With
\(\mathfrak D\mathbf x=\bar k-\bar m\) and
\(R_{\mathrm{ev}}=2\mathfrak D^\top\mathfrak D\), the restricted optimizer is
\(q^\star(\mathbf x,S)\propto\pi_Se^{-U_S(\mathbf x)/\mathrm T}\) with
\(U_S=V_S+\tfrac{\mathrm T}2|\mathbf x|^2+\tfrac{\mathrm T}2\mathbf x^\top R_{\mathrm{ev}}\mathbf x\),
and the physical posterior omits the last term. The kinetics are identity
mobility with diffusion covariance \(2\mathrm TI dt\) and one rate-one clock
at which the complete assignment is redrawn from the product of the receiver
posterior rows \(\beta_S(\mathbf x)\), self-proposals allowed; no generative
parameter is learned. The two perfect-matching assignments are
\(S^{(0)}=(1,0,3,2)\), which selects \(e_0,e_2\) from both ends, and
\(S^{(1)}=(3,2,1,0)\). For \(0<\epsilon\le1/100\) put
\begin{equation}
\mathrm T=\epsilon^8,\quad t_\epsilon=2\log(1/\epsilon),\quad
t^\circ_*=t_\epsilon+1,\quad \varrho_\circ=\tfrac1{100},\quad
\sigma_\circ=\varrho_\circ/\sqrt{48},\quad \mathbf x^\circ=(0,0,\mathbf 1_4).
\label{eq:rg-formation-cycle-regime}
\end{equation}
The preparation class consists of the laws
\(\frac1{16}\mathfrak p_\circ(\mathbf x)(1+\mathfrak z(\mathbf x)) d\mathbf x\)
on \((\mathbf x,S)\), where \(\mathfrak p_\circ\) is the equal mixture of
\(\mathcal N(\mathbf x^\circ,\sigma_\circ^2I)\) and
\(\mathcal N(-\mathbf x^\circ,\sigma_\circ^2I)\) and \(\mathfrak z\) is smooth,
invariant under the cycle automorphisms (with link signs transformed under
reflections), of \(\mathfrak p_\circ\)-mean zero, and with
\(\|\mathfrak z\|_\infty<1/2\).

The preparation names no block: its assignment is uniform and independent of
the configuration, reflections exchange its two centers, which carry opposite
loop orientations, and all initial edge mismatch second moments agree. Under
the source gating an edge that neither endpoint selects carries no
interaction energy, although every edge stays in the palettes of both its endpoints and every
attention weight stays positive.

\propositionheading{Endogenous pairing on a homogeneous four-cycle}{prop:rg-formation-endogenous-cycle}
In \Cref{def:rg-formation-cycle-experiment}, let \(0<\epsilon\le1/100\) and let
the initial law belong to the preparation class. The joint is normalized, the
scaled negative ELBO is
\(\Fenergy_{\mathrm T}(q)=\mathrm T\KL(q\Vert q^\star)+\mathrm B_{\mathrm T}
-\mathrm T\log Z_{\mathrm{opt}}\) with
\(\mathrm B_{\mathrm T}=\mathrm T[6\log(2\pi)+4\log2]\), and the recognition law
of the kinetics has nonincreasing \(\Fenergy_{\mathrm T}\) and is invariant
under the cycle automorphisms at every time. Put
\begin{equation}
p_\epsilon=73728 \epsilon^4t^\circ_*+4t^\circ_*e^{-1/\mathrm T}<\tfrac12,
\qquad
\eta_\epsilon=\tfrac92\varrho_\circ^2\epsilon^8+\frac{3}{4(1-e^{-8})}\epsilon^{32},
\qquad
\bar\eta_\epsilon=\frac{\eta_\epsilon+\mathcal S_2(p_\epsilon)}{1-p_\epsilon},
\label{eq:rg-formation-cycle-constants}
\end{equation}
with \(\mathcal S_2(p)=-p\log p-(1-p)\log(1-p)\). Then:

(a) \emph{Episode.} With probability at least
\(p_0=3e^{-1/4}/512\), the complete assignment equals one perfect matching
throughout \([0,t^\circ_*]\); at time zero the internal belief transported
relative entropy of each matched pair exceeds \(1/4\), so that the contrast is
entered during the episode and is not present at its start; on
\([t_\epsilon,t^\circ_*]\) both internal belief and model transported relative
entropies of each matched pair are at most \(\epsilon^4\); and on
\([1/4,t^\circ_*]\) every exterior belief relative entropy exceeds \(1/2\), the
model ones being larger. This is a coherence episode held by the two attended
edge costs: on the held matching sheet the normal restoring rates of the
kinetics are \(8+\mathrm T\) and \(4+5\mathrm T\), those of the physical law
\(8+\mathrm T\) and \(4+\mathrm T\), so at the residual resolution
\(\epsilon^2\) the least restoring rate times the squared resolution, divided
by the temperature, is \(4\epsilon^{-4}\), and the sheet is selected by the
posterior proposals on the event. The kinetics leave \(q^\star\) invariant,
and started from \(q^\star\) they realize the same event with probability at
most \(4e^{-1/(4\mathrm T)}\), so the episode exceeds its stationary
probability by at least \(p_0-4e^{-1/(4\mathrm T)}>p_0/2\).

(b) \emph{Parent channel.} For each matching \(M\), let \(Z_M\) be the span of
the three orthonormal directions
\((n^+_{e}-n^+_{f})/\sqrt2\), \(n^-_{e}\), \(n^-_{f}\), where \(e,f\) are its
two edges, \(n^\pm_e\) the unit normal covectors of
\cref{eq:rg-formation-cycle-normals}, and \(W_M=Z_M^\perp\). The channel
\(C(\mathbf x,S)=(S,P_{W_S}\mathbf x)\), with \(W_S=W_M\) on the two matchings
and \(W_S=\R^{12}\) otherwise, extended by the identity on observations and
evaluator samples, keeps every assignment, the loop
\(\sum_i\omega_i\), and a fast common-normal coordinate, and pushes the joint,
the physical posterior, the restricted optimizer, and the recognition law to
normalized parent laws.

(c) \emph{Discarded conditional free energy.} For every
\(t\in[t_\epsilon,t^\circ_*]\), the recognition law \(q_t\) of the episode
has mean discarded optimizing conditional relative entropy at most
\(\bar\eta_\epsilon\) and mean discarded physical conditional relative entropy,
with the evaluator samples retained, at most \(\bar\eta_\epsilon+\mathrm T\).

(d) \emph{Parent dynamics.} The inherited parent of \(C\), activated at
\(t_\epsilon\) from the complete recognition law of the episode prefix, has
retained path law on \([t_\epsilon,t^\circ_*]\), including the complete
assignment path, within total variation
\(2p_\epsilon/(1-p_\epsilon)+4\sqrt{8/3} e^{-1/\mathrm T}\) of the retained
native path conditioned on the same prefix.

The estimates hold under the stated source gating, random matching seed,
symmetric circulated preparation, temperature \(\epsilon^8\), Euclidean
kinetic metric, duration-one episode, and fixed generative parameters.
\status{ESTABLISHED}

The proof is in \cref{sec:rg-formation-cycle-proof}. On the center path
\cref{eq:rg-formation-cycle-center}, which holds \(S^{(0)}\) from
\(\mathbf x^\circ\), where every residual equals \(-1\), the selected
residuals decay as \(-e^{-(8+\mathrm T)t}\) while the exterior ones grow
toward \(-3/2\), so a source-cost gap greater than one develops by time
\(1/4\), and the posterior proposals then keep the matching except with
probability at most \(4t^\circ_*e^{-1/\mathrm T}\). Since the full
recognition law is invariant, the pairing is an event of sampled paths seeded
by the initial assignment, not a block chosen by the law.

The witness raises the question whether its mechanism, an edge that carries
energy only while an endpoint attends to it, produces a stationary law that
favors coherence beyond one level, or a second coalescence of the pairs it
forms. The next statement answers the stationary part negatively in the
class to which the four-cycle belongs, the family of
\Cref{def:rg-gated-translation-family} on an arbitrary finite graph; the
second coalescence is posed for models outside that family in
\Cref{open:rg-formation-general}. Under exclusive gating the source label
enters only the weight of a sheet and never its stiffness, so every sheet
restores its attended residual directions at rates of order one and leaves
every other direction, including every direction that translates one
coherent block relative to another, to the thermal reference at rates of
order \(\mathrm T\). This has a spectral and a stationary consequence.

\definitionheading{Gated translation family on a graph}{def:rg-gated-translation-family}
Let \((V,E)\) be a finite simple graph without isolated vertices, with
\(N=|V|\) agents and matching number \(\mathfrak n\), the largest number of
pairwise disjoint edges. Orient each edge \(e\in E\) once as \(e:j\to i\); the
interaction multigraph at \(c_*\) consists of the copies \(e\) and \(\bar e\)
of the edges of \((V,E)\), and the source set \(\mathcal N_i\) of receiver \(i\)
is its set of neighbors in \((V,E)\), with \(e(i,j)\) the edge joining \(i\) and
\(j\). The agents, links, laws, evaluator, and recognition coupling are those
of \Cref{def:rg-formation-cycle-experiment} on this graph: a configuration is
\(\mathbf x=(\bar k,\bar m,\omega)\in\R^{2N+|E|}\) with one link translation
\(\omega_e\) per edge, \(q_i^b=\mathcal N(\bar k_i,1)\),
\(q_i^m=\mathcal N(\bar m_i,1/2)\), and edge \(e:j\to i\) carries the residual
pair and the cost
\begin{equation}
\mathbf a_e\mathbf x=(r^b_e,r^m_e)
=(\bar k_i-\bar k_j-\omega_e,\ \bar m_i-\bar m_j-\omega_e),
\qquad
\mathcal D_e(\mathbf x)=\tfrac12|\mathbf a_e\mathbf x|^2,
\label{eq:rg-gated-residual-pair}
\end{equation}
so \(\mathbf a_e:\R^{2N+|E|}\to\R^2\) is linear and invariant under passive
frame changes as in \Cref{lem:rg-formation-carrier}, and \(\mathcal D_e\) is
the transported relative entropy \cref{eq:rg-formation-cycle-cost}. Each
receiver carries one source label of \Cref{def:obs-attention-label-model}
with source set \(\mathcal N_i\), so a complete assignment is
\(S\in\mathcal J=\prod_i\mathcal N_i\); the labels are exclusive, the prior is
the product \(\pi_S=\prod_i\pi_{iS_i}\) with every \(\pi_{ij}>0\) independent
of \(\mathrm T\), and given \((\mathbf x,S)\) independent binary records have
\(\Pr(E_i=1\mid\mathbf x,S)=\tfrac12\exp[-\mathcal D_{e(i,S_i)}(\mathbf x)/\mathrm T]\),
with \(E=\mathbf 1\) observed. The attended edge set, sheet energy, and sheet
stiffness of \(S\) are
\begin{equation}
\operatorname{att}(S)=\{e(i,S_i):i\in V\},\qquad
V_S=\sum_i\mathcal D_{e(i,S_i)}=\tfrac12\mathbf x^\top H_S\mathbf x,\qquad
H_S=\sum_i\mathbf a_{e(i,S_i)}^\top\mathbf a_{e(i,S_i)}.
\label{eq:rg-gated-sheet-stiffness}
\end{equation}
The configuration prior is \(\mathcal N(0,I_{2N+|E|})\), and with
\(R_{\mathrm{ev}}=2\mathfrak D^\top\mathfrak D\) the restricted optimizer and
the physical posterior are the laws
\begin{equation}
q_\Omega(\mathbf x,S)=\frac{\pi_S}{\sum_{S'}\pi_{S'}Z_{S'}}
\exp\Bigl[-\frac{V_S(\mathbf x)}{\mathrm T}-\tfrac12\mathbf x^\top\Omega\mathbf x\Bigr],
\qquad
Z_S=\int\exp\Bigl[-\frac{V_S(\mathbf x)}{\mathrm T}-\tfrac12\mathbf x^\top\Omega\mathbf x\Bigr]d\mathbf x,
\label{eq:rg-gated-sheet-law}
\end{equation}
with the reference precision \(\Omega=I+R_{\mathrm{ev}}\) for the optimizer
and \(\Omega=I\) for the physical posterior; \(Z_S\) is the sheet weight. For
either \(\Omega\) the kinetics have identity mobility, diffusion covariance
\(2\mathrm TI dt\), drift \(-\nabla_{\mathbf x}U_S\) with
\(U_S=V_S+\tfrac{\mathrm T}2\mathbf x^\top\Omega\mathbf x\), and one rate-one
clock at which the complete assignment is redrawn from the product of the
receiver posterior rows
\(\beta_i(j\mid\mathbf x)=\pi_{ij}e^{-\mathcal D_{e(i,j)}(\mathbf x)/\mathrm T}
\big/\allowbreak\sum_{j'\in\mathcal N_i}\pi_{ij'}e^{-\mathcal D_{e(i,j')}(\mathbf x)/\mathrm T}\),
self-proposals allowed; no generative parameter is learned. For an edge set
\(F\subseteq E\), \(\mathbf a_F\) stacks the maps \(\mathbf a_e\), \(e\in F\),
\(W_F=\ker\mathbf a_F\), \(\mathrm s_F\) is the least positive singular value
of \(\mathbf a_F\), and \(\operatorname{cc}(F)\) is the number of connected
components of the graph \((V,F)\), isolated vertices included. At resolution
\(\mathrm a>0\) the coherent edges of a configuration are
\(\operatorname{Coh}_{\mathrm a}(\mathbf x)=\{e\in E:|\mathbf a_e\mathbf x|\le\mathrm a\}\),
and \(\mathfrak P_{\mathrm a}(\mathbf x)\) is the partition of \(V\) into the
vertex sets of the components of \((V,\operatorname{Coh}_{\mathrm a}(\mathbf x))\),
the coherent blocks.

The four-cycle of \Cref{def:rg-formation-cycle-experiment} is the member with
\((V,E)\) the four-cycle and the uniform prior, and its kinetics are those with
\(\Omega=I+R_{\mathrm{ev}}\). A minimum edge cover of \((V,E)\) is an edge set
covering every vertex with the fewest edges; the proposition identifies these
covers with the spanning forests of stars with \(\mathfrak n\) components.

\propositionheading{Depth one under exclusive source gating}{prop:rg-gated-depth-one}
In \Cref{def:rg-gated-translation-family}, for either reference precision
\(\Omega\), let \(\mathrm k_S=\operatorname{rank}H_S\), let \(\mathrm h_S\) be
the least positive eigenvalue of \(H_S\), \(\mathrm h_*=\min_S\mathrm h_S\),
and \(\mathrm z_S=[\operatorname{pdet}H_S\cdot\det(\Omega|_{W_{\operatorname{att}(S)}})]^{-1/2}\),
where \(\operatorname{pdet}\) is the product of the positive eigenvalues and
\(\Omega|_W=\mathsf Q_W^\top\Omega\mathsf Q_W\) for a matrix \(\mathsf Q_W\)
whose columns form an orthonormal basis of \(W\); let \(\mathcal M\) be the set
of assignments of minimum rank \(\mathrm k_S\).

(a) \emph{Two restoring bands.} For every \(S\) and \(\mathrm T>0\), the
eigenvalues \(\lambda_1\le\cdots\le\lambda_{2N+|E|}\) of \(H_S+\mathrm T\Omega\)
satisfy \(\lambda_k\in[\mathrm T\|\Omega^{-1}\|^{-1},\mathrm T\|\Omega\|]\) for
\(k\le2N+|E|-\mathrm k_S\) and
\(\lambda_k\in[\mathrm h_S,\|H_S\|+\mathrm T\|\Omega\|]\) for
\(k>2N+|E|-\mathrm k_S\). The range of \(H_S\) is
\(W_{\operatorname{att}(S)}^\perp\), the span of the attended residual
covectors, and the prior enters \(q_\Omega\) only through the sheet weights
\(\pi_SZ_S\), never through \(H_S\).

(b) \emph{Gating wells.} With
\[
V_\pi(\mathbf x)=-\mathrm T\sum_i\log\sum_{j\in\mathcal N_i}\pi_{ij}e^{-\mathcal D_{e(i,j)}(\mathbf x)/\mathrm T},
\qquad
V_{\min}(\mathbf x)=\sum_i\min_{j\in\mathcal N_i}\mathcal D_{e(i,j)}(\mathbf x),
\]
one has
\begin{equation}
0\le V_\pi(\mathbf x)-V_{\min}(\mathbf x)
\le\mathrm T\sum_i\log\frac1{\min_{j\in\mathcal N_i}\pi_{ij}}
\qquad\text{for every }\mathbf x,
\label{eq:rg-gated-well-depth}
\end{equation}
the configuration marginal of \(q_\Omega\) has density proportional to
\(\exp[-V_\pi(\mathbf x)/\mathrm T-\tfrac12\mathbf x^\top\Omega\mathbf x]\),
and if every edge of an edge set \(F\) covering \(V\) has zero residual at
\(\mathbf x\), then \(V_{\min}\) vanishes at \(\mathbf x\) and at every
configuration obtained from it by translating the belief and model means of
the agents of one component of \((V,F)\) by a common amount with the links
unchanged.

(c) \emph{Rank, minimum edge covers, and sheet weights.} For every
\(F\subseteq E\), \(\operatorname{rank}\mathbf a_F=|F|+N-\operatorname{cc}(F)\).
Hence \(\mathrm k_S=|\operatorname{att}(S)|+N-\operatorname{cc}(\operatorname{att}(S))\ge2(N-\mathfrak n)\),
with equality exactly when \(\operatorname{att}(S)\) is a minimum edge cover
of \((V,E)\); the minimum edge covers are the spanning forests of stars with
\(\mathfrak n\) components, every one of them is attended by some
assignment, and an edge outside \(\operatorname{att}(S)\) joining two of its
components raises the rank by two. The sheet weight satisfies
\begin{equation}
\Bigl(1+\frac{\mathrm T\|\Omega\|}{\mathrm h_S}\Bigr)^{-\mathrm k_S/2}
\le\frac{Z_S}{(2\pi)^{(2N+|E|)/2}\mathrm z_S\mathrm T^{\mathrm k_S/2}}\le1,
\label{eq:rg-gated-sheet-weight}
\end{equation}
and consequently, for \(0<\mathrm T\le1\) and every integer \(m\ge1\),
\begin{equation}
q_\Omega\{S:\operatorname{cc}(\operatorname{att}(S))\le\mathfrak n-m\}\le\mathrm K_m\mathrm T^m,
\qquad
\mathrm K_m=\Bigl(1+\frac{\|\Omega\|}{\mathrm h_*}\Bigr)^{N-\mathfrak n}
\frac{\sum_{S:\operatorname{cc}(\operatorname{att}(S))\le\mathfrak n-m}\pi_S\mathrm z_S}
{\sum_{S\in\mathcal M}\pi_S\mathrm z_S}.
\label{eq:rg-gated-merge-cost}
\end{equation}

(d) \emph{Stationary block law.} Let \(\mathrm s_*\) be the least, over all
\(S\) and all edges \(e\notin\operatorname{att}(S)\) joining two components of
\(\operatorname{att}(S)\), of the least positive singular value of
\(\mathbf a_eP_{W_{\operatorname{att}(S)}}\), where \(P_W\) is the orthogonal
projection onto \(W\). When no such pair \((S,e)\) exists, put
\(\mathrm s_*=+\infty\) and \(1/\mathrm s_*^2=0\). For \(\mathrm a>0\) and
\(0<\mathrm T\le1\), the
probability under \(q_\Omega\) that \(\mathfrak P_{\mathrm a}(\mathbf x)\) is
not the partition of \(V\) into the components of a minimum edge cover is at
most
\begin{equation}
\mathrm K_1\mathrm T+\frac{|E|\|\Omega\|}{2\mathrm s_*^2}\mathrm a^2+4|E|e^{-\mathrm a^2/(4\mathrm T)}.
\label{eq:rg-gated-block-law}
\end{equation}
\status{ESTABLISHED}

The proof is in \cref{sec:rg-formation-depth-one-proof}. In the four-cycle,
\(\mathfrak n=2\) and the minimum edge covers are the two perfect matchings,
so (c) gives them all but \(O(\mathrm T)\) of the stationary mass of either
law and (d) makes coherence of an unattended bridge at resolution
\(\mathrm a\) cost \(O(\mathrm a^2)\) per bridge. A second coalescence that
is held by an attended bridge at stationarity with probability bounded below
as \(\mathrm T\to0\) therefore needs a model outside the family.

\openproblemheading{Repeated endogenous formation}{open:rg-formation-general}
Construct, at \(c_*\), a second endogenous coalescence with inherited source
interactions and retained marks, in a family where a matching or partition
forms from initially nonmatching source assignments, generative parameters are
learned, and the episode length is not fixed in advance. What holds is
\Cref{prop:rg-formation-endogenous-cycle}, one duration-one episode seeded by
the initial assignment, \Cref{thm:finite-depth-adaptive-formation}, whose
hierarchy is supplied, and \Cref{prop:rg-gated-depth-one}, by which in the
gated translation family of \Cref{def:rg-gated-translation-family} the
stationary coherent-block partition is, up to the error of
\cref{eq:rg-gated-block-law}, the partition into the components of a minimum
edge cover, and every further merge at stationarity costs a factor
\(\mathrm T\). The question includes arbitrary families and initial
laws, mixed-level source support, permanent membership under later forcing,
and control as the depth tends to infinity; the intrinsic selector and the
hierarchy question are
\Cref{open:cg-natural-selector,open:full-graph-hierarchy}. It is posed for
models at \(c_*\) with source-gated attention on a finite connected graph
whose energies, precisions, prior ratios, hyperparameters, and mobilities do
not depend on the temperature \(\mathrm T\), as \(\mathrm T\to0\), outside the
family of \Cref{def:rg-gated-translation-family}. A second coalescence in
such a model is declared by a handoff time, deterministic or a stopping time;
a first-stage event of the history up to the handoff, with its block
partition; a second-stage event, with windows measured from the handoff,
whose block is a union of at least two first-stage blocks; one statistic of
the handoff state, namely the probability of the second-stage event, the
probability that it has occurred by declared times, or the probability of a
declared block; and a handoff class, fixed before the comparison, that
contains every first-stage handoff and, with probability bounded below, the
handoffs of histories without first-stage formation. The problem is settled
by such a coalescence, with learned generative parameters and from initially
nonmatching assignments, whose second-stage event does not occur with
probability tending to one uniformly over the handoff class, whose statistic
after first-stage formation exceeds its value after a counterfactual handoff
in the class by a margin bounded below, whose block is held at a declared
resolution by a restoring interaction whose rate times the squared
resolution, divided by \(\mathrm T\), tends to infinity on the sheets the
kinetics select, on horizons short compared with the stationary recurrence
time of the event measured by its rate of entrance from a set of stationary
mass bounded below, and with parents whose conditional and dynamical errors
vanish separately; or by a proof that in this class every coalescence of
first-stage blocks on such horizons fails the margin or the restoring
clause or has a second-stage event occurring with probability tending to one
uniformly over the handoff class.

\section*{Physicist's summary}

In the triangle chain a supplied hierarchy of stiffnesses
\(\mathrm J_1\gg\cdots\gg\mathrm J_L=1\) makes the prepared displacement relax
like overdamped springs of very different stiffness: the stiffest bonds
equilibrate first and glue their endpoints, and each softer level glues the
clusters already formed (\Cref{thm:finite-depth-adaptive-formation}). Every
parent keeps every loop and source assignment, starts from the complete
current law, and learns its own parameter; the price of each reduction
vanishes in an extreme regime (temperature \(\epsilon^{86}\) at depth two) in
which nothing adaptive decides the outcome. In the homogeneous four-cycle, without
a supplied hierarchy, an edge carries energy only when an endpoint attends to
it; attended edges contract, the prepared circulation stretches the unattended
ones, and the posterior keeps attending to the contracted edges, so a pair
forms with probability bounded away from zero
(\Cref{prop:rg-formation-endogenous-cycle}). The same gating
is why the story stops there in that family
(\Cref{prop:rg-gated-depth-one}): an unattended edge between two formed
blocks feels only the thermal reference, so at stationarity the blocks are
the components of a minimum edge cover and each further merge costs a factor
\(\mathrm T\).
